%% versione preprint su arxiv 8 settembre
\documentclass[a4paper,10pt]{article}
\usepackage{amsmath}
\usepackage{amsfonts}
\usepackage{amssymb}

\newtheorem{theorem}{Theorem}
\newtheorem{acknowledgement}[theorem]{Acknowledgement}
\newtheorem{corollary}[theorem]{Corollary}
\newtheorem{criterion}[theorem]{Criterion}
\newtheorem{definition}[theorem]{Definition}
\newtheorem{exercise}[theorem]{Exercise}
\newtheorem{lemma}[theorem]{Lemma}
\newtheorem{notation}[theorem]{Notation}
\newtheorem{problem}[theorem]{Problem}
\newtheorem{proposition}[theorem]{Proposition}
\newtheorem{remark}[theorem]{Remark}
\newtheorem{solution}[theorem]{Solution}
\newtheorem{summary}[theorem]{Summary}
\newtheorem{hypothesis}[theorem]{Hypothesis}
\newenvironment{proof}[1][Proof]{\textbf{#1.} }{\ \rule{0.5em}{0.5em}}




\begin{document}


\title{
%Well posedness of transport equations with H\"{o}lder vector fields by stochastic perturbation
Well-posedness of the transport equation by stochastic
perturbation }
\author{F. Flandoli$^{1}$,
M. Gubinelli$^{2}$, E. Priola$^{3}$\\{\small {(1) Dipartimento di
Matematica Applicata ``U. Dini'', Universit\`a di Pisa,
Italia}}\\{\small { (2) CEREMADE  (UMR 7534), Universit\'e Paris
Dauphine, France}}\\{\small { (3) Dipartimento di Matematica,
Universit\`a di Torino, Italia }}} \maketitle
\begin{abstract}
We consider  the linear transport equation with a globally
H\"{o}lder continuous and bounded vector field. While this
deterministic PDE may not be well-posed, we prove that a
multiplicative stochastic perturbation of Brownian type is enough
to render the equation well-posed. This seems to be the first
explicit example of partial differential equation that become
well-posed under the influece of noise. The key tool is a
differentiable stochastic flow constructed and analysed by means
of a special transformation of the drift of Tanaka type.
\end{abstract}





\section{Introduction}

The transport equation in $[0,T]\times\mathbb{R}^{d}$%
\begin{align*}
\partial_{t}u(t,x)+b(t,x)\cdot\nabla u(t,x)  &  =0,\quad(t,x)\in
\lbrack0,T]\times\mathbb{R}^{d},\\
u(0,x)  &  =u_{0}(x),\quad x\in\mathbb{R}^{d}.
\end{align*}
driven by a vector field $b:[0,T]\times\mathbb{R}^{d}\rightarrow\mathbb{R}%
^{d}$ which is Lipschitz continuous in $x$, uniformly in $t$, can
be solved in several classical ways and in different function
spaces. A remarkable extension has been obtained by R.J. Di Perna
and P.L. Lions~\cite{DiPernaLions} who proved, under the
assumption $b\in L^{1}( 0,T;W_{loc}^{1,1}(
\mathbb{R}^{d},\mathbb{R}^{d}))  $ with a linear growth condition
and ${\mathrm{div}} b \in L^1(0, T; L^{\infty}(\mathbb {R}^d))$,
that a unique
 $L^{\infty}( 0,T;L^{\infty }( \mathbb{R}^{d}) ) $ solution exists,
weak-$*$
continuous in time, for a given $u_{0}\in L^{\infty}(  \mathbb{R}%
^{d})  $. Moreover, a generalized notion of flow is introduced and its
existence and uniqueness is proved under the assumption that the negative part
of ${\mathrm{div}\,}b$ is $L^{1}(  0,T;L^{\infty}(  \mathbb{R}%
^{d}) )  $. L. Ambrosio \cite{Ambrosio} extended further the
theory including in particular the case of vector fields $b\in L^{1}(
0,T;BV_{loc}(  \mathbb{R}^{d},\mathbb{R}^{d}) )  $ with
similar conditions on ${\mathrm{div}\,}b$. The literature following
\cite{DiPernaLions} is wide, see a partial review in \cite{AmbrCrippa}. See
also a generalization to transport-diffusion equations and the associated
stochastic differential equations by Figalli \cite{Figalli} (in a direction
different from the one of the present paper).

In all these results some regularity assumption on $\nabla b$ is
imposed, at least $L^{1}$ or a signed finite measure. If we assume
only that $b$ is H\"{o}lder continuous, couterexamples to
 well-posedness can be constructed.

The aim of this paper is to show that suitable random
perturbations of the transport equation improve the theory. We
prove existence and uniqueness of
$L^{\infty}$-solutions when $b$ is continuous, bounded, globally ${\alpha}%
$-H\"{o}lder continuous in space, ${\alpha}\in(0,1)$, uniformly in
time, and satisfy a suitable  integrability condition on
${\mathrm{div}\,}b$. This result is obtained at the price of
introducing a multiplicative white noise in the PDE. Precisely, we
consider the equation
\begin{equation}
%\left\{
\begin{aligned}
 \partial_{t}u(  t,x)  &+(  b(  t,x)  \cdot\nabla
u(  t,x) )  dt+\sum_{i=1}^{d}e_{i}\cdot\nabla u( t,x)  \circ
dW_{t}^{i}=0,\\  u(  0,x) & =u_{0}( x),\;\; x \in \mathbb{R}^d,
\end{aligned}
%\right.
\label{SPDE}%
\end{equation}
where $e_{1},...,e_{d}$ is the canonical basis of ${\mathbb{R}}^{d}$ and
$W_{t}=(W_{t}^{1},...,W_{t}^{d})$ is a standard Brownian motion in
${\mathbb{R}}^{d}$. The stochastic integration is understood in the
Stratonovich sense.

This equation  seems to be the first concrete example  of partial
differential equation that is not well-posed without noise, but is
well-posed with a suitable noise.
%is one of the few known examples (perhaps the first concrete one)
%of partial differential equation that is not wellposed without
% noise, but is wellposed with a suitable noise.
This phenomenon is well known for ordinary differential equations
but it is one of the more interesting direction of investigations
in SPDEs, where people have in mind the hope to prove
well-posedness for some relevant equation of fluid dynamics.

The well-posedness of equation~(\ref{SPDE}) relies on two facts.
The associated stochastic ordinary differential equation (the
stochastic characteristics)
\begin{equation}
dX_{t}^{x}=b\left(  t,X_{t}^{x}\right)  dt+dW_{t},\quad t\geq0,\quad X_{0}%
^{x}=x\label{SDE intro}%
\end{equation}
under the assumption $b(t,\cdot)\in C_{b}^{{\alpha}}({\mathbb{R}}%
^{d};{\mathbb{R}}^{d})$, uniformly in time, generates a stochastic flow of
diffeomorphisms. Ultimately, the reason is that integrals of the form
\begin{equation}
\int_{0}^{t}f\left(  s,X_{s}^x\right)  ds\label{special integrals}%
\end{equation}
 are more regular in $x$ in the case of diffusion processes $X_{s}$ than
in the case of deterministic solutions to ODEs. The existence
of a stochastic flow of diffeomorphisms implies the existence of
 weak solutions to equation~(\ref{SPDE}) in $L^{\infty}$.

The uniqueness of such weak solutions is more difficult and, even
under the existence of a stochastic flow of diffeomorphisms, it
must be related to some form of commutator lemma, as in the
deterministic case (see \cite{DiPernaLions}, \cite{Ambrosio}),
since we need to perform differential computations on
regularizations of $L^{\infty}$-solutions. We use a weak form for
the commutator composed with the flow, which does not require weak
differentiability of $b(  t,\cdot)  $ but only requires weak
differentiability of the Jacobian of the flow itself. The Jacobian
of the flow
is given by $\exp(  \int_{0}^{t}{\mathrm{div}\,}b(  X_{s}%
^{x})  ds)  $. In the deterministic case there is no hope to
differentiate in $x$ without further differentiability assumptions
on ${\mathrm{div}\,}b$, but in the stochastic case we use again
the improved regularity of integrals of the form (\ref{special
integrals}). We have only to impose an integrability condition on
${\mathrm{div}\,}b$, which turns out to be different from
 the one required in the deterministic setting \cite{DiPernaLions},
\cite{Ambrosio}.

The generalization to nonlinear transport equations, where $b$
depends on $u$ itself, would be a major next step for applications
to fluid dynamics but it turns out to be a  difficult problem.
What looks very difficult is to generalize to an adapted random
field $b(  t,x,\omega)  $. Indeed, it is very easy to
produce examples, both for the linear SPDE (\ref{SPDE}) and
for a stochastic version of Euler equation which show that the particular noise we use does not have any regularizing effect.  Thus new ideas are
necessary to approach nonlinear problems.

\paragraph{Plan.}

In Sect.~\ref{sec:ito} we prove the existence of a global
stochastic flow associated to eq.~(\ref{SDE intro}) and its
differentiability properties. In Sect.\ref{sec:transp} we use the
stochastic flow to prove existence  of weak solutions to the
SPDE~(\ref{SPDE}).  Sect. 4 and 5 are devoted to prove uniqueness
for solutions to the SPDE~(\ref{SPDE}). Finally, in
 Sect.~\ref{sec:counter} we point out the basic difficulty present when trying to
 use the same approach for non-linear transport equations arising
 from fluid dynamics.

\paragraph{Notations.}

Let $T>0$ be fixed. For ${\alpha}\in(0,1)$ define the space $C_{b}^{\alpha
}([0,T]\times{\mathbb{R}}^{d};{\mathbb{R}})$ as the set of all bounded
continuous functions $f:[0,T]\times{\mathbb{R}}^{d}\rightarrow{\mathbb{R}}$
for which
\[
\lbrack f]_{\alpha,T}:=\sup_{t\in\lbrack0,T]}\sup_{x\neq y\in{\mathbb{R}}^{d}%
}\frac{|f(t,x)-f(t,y)|}{|x-y|^{\alpha}}<\infty
\]
($|\cdot|$ denotes the euclidean norm in ${\mathbb{R}}^{d}$ for every $d$, if
no confusion may arise). This is a Banach space with respect to the usual norm
$\Vert f\Vert_{{\alpha},T}=\Vert f\Vert_{0}+[f]_{\alpha,T}$ where $\Vert
f\Vert_{0}=\sup_{(t,x)\in\lbrack0,T]\times{\mathbb{R}}^{d}}|f(t,x)|$.

We write
$C_{b}^{\alpha}([0,T]\times{\mathbb{R}}^{d};{\mathbb{R}}^{d})$ or
simply $C_{b,T}^{{\alpha}}$ for the space of all vector fields
$f:[0,T]\times {\mathbb{R}}^{d}\rightarrow{\mathbb{R}}^{d}$ having
all components in
$C_{b}^{\alpha}([0,T]\times{\mathbb{R}}^{d};{\mathbb{R}})$.

Moreover $f\in C_{b}^{n+{\alpha}}([0,T]\times{\mathbb{R}}^{d};{\mathbb{R}}%
^{d})$, briefly $f\in C_{b,T}^{n+{\alpha}}$, if for all
$k=0,\dots,n$, the spatial Fr\'echet derivatives $D^{k}f$ are all
bounded and $\theta$-H\"{o}lder continuous. Define the
corresponding norm as
\[
\Vert f\Vert_{n+\alpha,T}=\Vert f\Vert_{0}+\sum_{k=1}^{n}\Vert D^{k}f\Vert
_{0}+[D^{n}f]_{\alpha,T}.
\]
The spaces $C_{b}^{n+{\alpha}}=C_{b}^{n+{\alpha}}({\mathbb{R}}^{d}%
,{\mathbb{R}}^{d})$ are defined as before but only involve functions
$f:{\mathbb{R}}^{d}\rightarrow{\mathbb{R}}^{d}$ which do not depend on time.

We say  that $f: \mathbb{R} ^d \to \mathbb{R} ^d$ belongs to
$C^{k, \alpha}$, $k \in \mathbb N$, $\alpha \in (0,1)$, if $f$ is
continuous on $\mathbb{R} ^d$, $k$-times differentiable with all
continuous derivatives and  the derivatives of order $k$ are
locally $\alpha$-H\"older continuous.

Finally, $C^0_0(\mathbb{R} ^d)$ denotes the space of all real
continuous functions defined on $\mathbb{R} ^d$, having compact
support and $ C^{\infty}_0(\mathbb{R} ^d)$ its subspace consisting
of infinitely differentiable functions.



\vskip 1mm Throughout the paper we will assume a stochastic basis
with a $d$-dimensional Brownian motion $\left(  \Omega,\left(
\mathcal{F}{}_{t}\right) ,{}\mathcal{F},P,\left(  W_{t}\right)
\right)  $ to be given. Denote by $\mathcal{F}_{s,t}$ the
 completed
  $\sigma$-algebra generated by $W_{u}-W_{r}$, $s\leq r\leq u\leq
 t$,
 for each $0\le s<t$, and set $\mathcal{F}_{s,s} = \{ \varnothing, \Omega
  \}$.

We will consider the following assumptions.

\begin{hypothesis}
\label{hy1}
There exists ${\alpha}\in(0,1)$ such that $b\in C_{b,T}^{{\alpha}}$.
\end{hypothesis}

\begin{hypothesis}
\label{hy2}
There exists $p,  q \in (2, +\infty) $ such that
 \begin{equation} \label{uni}
{\mathrm{div}}b \in L^{q} (0,T; L^p(\mathbb{R} ^d))  \;\;\;\;
\text{with} \;\; \frac{d}{p} + \frac{2}{q} < 1
 \end{equation}
(where ${\mathrm{div}\,}b(t,\cdot)$ is understood in distributional sense).
\end{hypothesis}


 \begin{remark} In \cite{Kry-Ro} condition \eqref{uni}
 is imposed  locally on $b$ in order to prove existence and
 uniqueness of strong solution to \eqref{SDE intro}.
  This hypothesis     is satisfied
if $d \ge 2$ and
$$
{\mathrm{div}} b \in L^{\infty} (0,T; L^p(\mathbb{R} ^d))
\;\;\;\; \text{for some} \;\; p
>d.
$$
\end{remark}
For a given function $f$ depending on $t\in\lbrack0,T]$ and $x\in{\mathbb{R}%
}^{d}$, we will also adopt the  notation
\[
f_{t}(x)=f(t,x).
\]


\section{Differentiable stochastic flow for an SDE with $C_{b,T}^{{\alpha}}$-drift}

\label{sec:ito}

\textit{In this section we assume that $b$ satisfies  Hypothesis
\ref{hy1}. }

{Given }$s\in\left[  0,T\right]  $ and $x\in{\mathbb{R}}^{d}$, consider the
stochastic differential equation in ${\mathbb{R}}^{d}$ :
\begin{equation}
dX_{t}=b\left(  t,X_{t}\right)  dt+dW_{t},\quad t\in\left[  s,T\right]  ,\quad
X_{s}=x. \label{SDE}%
\end{equation}
A classical fact is the existence and uniqueness of a weak
solution, obtained for instance by Girsanov transform.
Yu.~Veretennikov \cite{V} prove that boundedness of $b$
(independent on $t$) is enough to have path-wise uniqueness and
existence of strong probabilistic solutions. For related work see
also the recent paper~\cite{Kry-Ro} by N.V.~Krylov and
M.~R\"ockner. Recently, A.M.~Davie has proved that for a.e.
Brownian path $W$ it is possible to prove uniqueness for
equation~(\ref{SDE})
between all deterministic continuous functions which solve equation~(\ref{SDE}%
) with respect to the choosen $W$ (for any initial condition and
any bounded measurable and time dependend drift $b$). The analysis
of the associated flow is however still missing.

\vskip 3mm

Our key result is the existence of a \emph{differentiable} stochastic flow
$(x,s,t)\mapsto\varphi_{s,t}(x)$ for equation~(\ref{SDE}). Recall the relevant
definition from H. Kunita~\cite{K}:

\begin{definition}
A \emph{stochastic flow of diffeomorphisms} (resp. \emph{of class}
$C^{1,\alpha}$) on $\left(  \Omega,\left(  \mathcal{F}{}_{t}\right)
,{}\mathcal{F},P,\left(  W_{t}\right)  \right)  $ associated to equation
(\ref{SDE}) is a map $(s,t,x,\omega)\mapsto\phi_{s,t}(x)\left(  \omega\right)
$, defined for $0\leq s\leq t\leq T$, $x\in{\mathbb{R}}^{d}$, $\omega\in
\Omega$ with values in ${\mathbb{R}}^{d}$, such that

\begin{itemize}
\item [(a)]given any $s\in\left[  0,T\right]  $, $x\in{\mathbb{R}}^{d}$, the
process $X^{s,x}=\left(  X_{t}^{s,x}\left(  \omega\right)  ,t\in\left[
s,T\right]  ,\omega\in\Omega\right)  $ defined as $X_{t}^{s,x}=\phi_{s,t}(x)$
is a continuous $\mathcal{F}_{s,t}$-measurable solution of equation (\ref{SDE})

\item[(b)] $P$-a.s., $\phi_{s,t}$ is a diffeomorphism, for all
$0\leq s\leq t\leq T$, and the functions $\phi_{s,t}(x)$,
$\phi_{s,t}^{-1}(x)$, $D\phi_{s,t}(x)$, $D\phi_{s,t}^{-1}(x)$ are
continuous in $(s,t,x)$ (resp. of class $C^{1,\alpha}$ in $x$
uniformly in $(s,t)$),

\item[(c)] $P$-a.s., $\phi_{s,t}(x)=\phi_{s,u}(\phi_{u,t}(x))$ for all $0\leq
s\leq u\leq t\leq T$ and $x\in{\mathbb{R}}^{d}$ and $\phi_{s,s}(x)=x$.
\end{itemize}
\end{definition}

 As already mentioned
the main ingredient  to obtain the regularity of the flow is the
observation that the time integral
$\int_{0}^{t}b\left(s,X_{s}^x\right)  ds$ has richer regularity
properties than expected only on the basis of the regularity of
$b$; such properties, coming from the interaction between $b$ and
the local behavior of $W_{t}$, and ultimately related to the local
time (or occupation measure) of the Brownian motion $W_{t}$, may
be revealed by the Tanaka approach to the definition of local
time: the integral $\int_{0}^{t}b\left(  s,X_{s}^x\right)  ds$ may
be considered as the corrector of an It\^{o} formula, and the
other terms of the formula have better properties than
$\int_{0}^{t}b\left( s,X_{s}^x\right)  ds$.



\medskip
Let $\lambda>0$ be fixed. Let us extend $b$ to the whole
$[0,\infty )\times{\mathbb{R}}^{d}$ by setting
 \begin{equation} \label{bb}
 b(t,x)=b(T,x),\;\;\;
t\geq T,\;\;  x\in{\mathbb{R}}^{d}.
 \end{equation} Given a function
$f:[0,+\infty)\times{\mathbb{R}} ^{d}\rightarrow{\mathbb{R}}^{d}$
consider the following parabolic
 system (collecting $d$ independent equations):
\begin{equation}
\partial_{t}u_{\lambda}+L^{b}u_{\lambda}-\lambda u_{\lambda}=f,\;\;\;(t,x)\in
(0,\infty)\times{\mathbb{R}}^{d} \label{eq:parabolic},
\end{equation}
where
\begin{equation} \label{fr}
L^{b}u=\frac{1}{2}\Delta u+b\cdot\nabla u
 \end{equation} and
 $u:[0,+\infty)\times{\mathbb{R}}
 ^{d}\rightarrow{\mathbb{R}}^{d}$ (clearly \eqref{fr} has to be interpreted
 componentwise).



\vskip 1mm The next  result deals with Schauder estimates and is
 known even in a more general form
 (see  \cite{KP} and the references therein). Schauder estimates
   of the type \eqref{stima} were first proved in  \cite{Br}.
  We will only sketch a proof for the
 sake of completeness.

\begin{lemma}
\label{schaud} Let $\lambda>0$, $b \in C^{\alpha}_{b,T}$, $\alpha
\in (0,1),$ and $f:[0,+\infty)\times{\mathbb{R}}
^{d}\rightarrow{\mathbb{R}}^{d}$ be a continuous and bounded such
that $f(t,\cdot)\in
C_{b}^{\alpha}({\mathbb{R}}^{d};{\mathbb{R}}^{d})$, $t\geq0$. Then
there exists a unique bounded classical solution $u = u_{\lambda}$
to equation
(\ref{eq:parabolic}). The functions $u:[0,\infty)\times{\mathbb{R}}%
^{d}\rightarrow{\mathbb{R}^d}$ and $\partial_{t}u$ are continuous.
Moreover there exists $C>0$ (independent on $u$ and $f$) such that
\begin{equation}
\sup_{t\geq0}\Vert u(t,\cdot)\Vert_{C_{b}^{2+\alpha}}\leq
C\sup_{t\geq0}\Vert f(t,\cdot)\Vert_{C_{b}^{\alpha}} \label{stima}
\end{equation}
\end{lemma}

\begin{proof}
Uniqueness follows by the maximum principle (see \cite{Kr}). Let us consider existence.

If
   $L^{b}=\frac{1}{2}\triangle$, then the result follows
by using the explicit formula
\begin{equation}
u(t,x)=\int_{t}^{+\infty}e^{-\lambda(r-t)}P_{r-t}f(r,\cdot)(x)dr=\int
_{0}^{+\infty}e^{-\lambda s}P_{s}f(t+s,\cdot)(x)ds \label{heat}%
\end{equation}
(where $(P_{t})$ denotes the heat semigroup) and well known
estimates on the spatial derivatives of $P_t g$ when $g \in
C^{\alpha}_b$ and $t>0$. Let
$L^{b}=\frac{1}{2}\triangle+b\cdot\nabla$. Thanks to the
boundedness of $b$ we get easily a-priori estimates for equation
(\ref{eq:parabolic}) (assuming that there exists a
 classical bounded solution $u$  such that
$\sup_{t\geq0}\Vert u(t,\cdot)\Vert_{C_{b}^{2+\alpha}}<\infty$).
Then a standard continuity method (see for instance \cite{Kr})
allows to get the existence of the solution which verifies
equation (\ref{eq:parabolic}). The proof is complete.
\end{proof}

\vskip 1mm  We also need the following simple lemma.


\begin{lemma}
\label{grad} Under the assumptions of Lemma \ref{schaud}
 let $u_{\lambda}$ be the solution to  (\ref{eq:parabolic}%
). Then
\[
\sup_{t\geq0}\Vert \nabla u_{\lambda}(t,\cdot)\Vert_{0}\rightarrow
0,\;\;\;as\;\;\lambda\rightarrow+\infty.
\]
\end{lemma}

\begin{proof}
We write $\partial_{t}u_{\lambda}+\frac{1}{2}\triangle u_{\lambda}-\lambda
u_{\lambda}=f-b\cdot\nabla u_{\lambda}.$

Using the well known estimate for the heat semigroup $\Vert\nabla P_{t}%
g\Vert_{0}\leq\frac{C}{t^{\frac{1}{2}}}\Vert g\Vert_{0},$ $g\in C_{b}%
({\mathbb{R}}^{d}),$ $t>0,$ and, differentiating in equation
(\ref{eq:parabolic}), we get easily, for any $\lambda>0$,
\[
\Vert\nabla u_{\lambda}\Vert_{0}\leq\frac{c}{\lambda^{\frac{1}{2}}}(\Vert
f\Vert_{0}+\Vert b\cdot\nabla u_{\lambda}\Vert_{0})\leq\frac{c}{\lambda
^{\frac{1}{2}}}\Vert f\Vert_{0}+\frac{c}{\lambda^{\frac{1}{2}}}\Vert
b\Vert_{0}\Vert\nabla u_{\lambda}\Vert_{0},
\]
where the suprema are taken on $[0,\infty)\times{\mathbb{R}}^{d}$. Considering
$\lambda>c^{2}\Vert b\Vert_{0}^{2}$, we get
\[
(1-\frac{c}{\lambda^{\frac{1}{2}}}\Vert b\Vert_{0})\,\Vert\nabla u_{\lambda
}\Vert_{0}\leq\frac{c}{\lambda^{\frac{1}{2}}}\Vert f\Vert_{0}%
\]
and the assertion follows. The proof is complete.
\end{proof}


\begin{theorem}
\label{th:flow1} Assume $b\in C_{b,T}^{{\alpha}}$.
 Then we have the following facts.

\begin{itemize}
\item [(i)](pathwise uniqueness) For every $s\in\left[  0,T\right]
$, $x\in{\mathbb{R}}^{d}$, the stochastic equation (\ref{SDE}) has
a unique continuous adapted solution $X^{s,x}=\left(
X_{t}^{s,x}\big( \omega\right) ,t\in\left[ s,T\right],$
$\omega\in\Omega\big)  $.

\item[(ii)] (differentiable flow) There exists a stochastic flow $\phi_{s,t}$
of diffeomorphisms for equation (\ref{SDE}). The flow is also of class
$C^{1+{\alpha}^{\prime}}$ for any ${\alpha}^{\prime}<{\alpha}$.

\item[(iii)] (stability) Let $(b^{n})\subset C_{b,T}^{{\alpha}}$
be a sequence of vector fields and $\phi^{n}$ the corresponding
stochastic flows. If $b^{n}\rightarrow b$ in $C_{b,T}^{\alpha}$,
then, for any $p\geq 1$,
\begin{equation}
\label{stability1}
\lim_{n\rightarrow\infty}\sup_{x\in{\mathbb{R}}^{d}} \sup_{0\leq s\leq
T} E[ \sup_{u \in [s,T]} |\phi_{s,u}^{n}(x)-\phi_{s,u}(x)|^{p}]=0.
\end{equation}%
\begin{equation}
\label{stability2}
\lim_{n\rightarrow\infty}\sup_{x\in{\mathbb{R}}^{d}}\sup_{0\leq s\leq
T} E[  \sup_{u \in [s,T]} \Vert D\phi_{s,u}^{n}(x)-D\phi_{s,u}(x)\Vert^{p}]=0.
\end{equation}
(here $\Vert\cdot\Vert$ denotes the Hilbert-Schmidt norm).
\end{itemize}
\end{theorem}

\begin{proof}
\textbf{Step 1} (auxiliary parabolic systems). For $\lambda>0$ consider the
(vector valued) solution $\psi$ to the parabolic system
 \begin{equation} \label{pde}
\partial_{t}\psi_{\lambda}+L^{b}\psi_{\lambda}-\lambda\psi_{\lambda
}=-b,\;\;\;(t,x)\in(0,\infty)\times{\mathbb{R}}^{d}%
\end{equation}
provided by an application of Lemma~\ref{schaud} (recall
\eqref{bb}). Define
\[
\Psi_{\lambda}(t,x)=x+\psi_{\lambda}(t,x).
\]
%Notice that
%\begin{equation}
%\partial_{t}\Psi_{\lambda}+L^{b}\Psi_{\lambda}=\lambda\psi_{\lambda}.
%\label{PDE capital psi}%
%\end{equation}

\begin{lemma} \label{diff} For  $\lambda $ large enough, such that
 $\sup_{t\geq0}\Vert \nabla \psi_{\lambda}(t,\cdot)\Vert_{0} <1$ (see
 Lemma \ref{grad}),
 the following statements hold:

{{\vskip 2mm \noindent }} (i) Uniformly in $t \in [0,+\infty)$,
$\Psi_{\lambda}$ has bounded first and second spatial
 derivatives and moreover  the second (Fr\'echet) derivative $D^2_x
\Psi_{\lambda}$ is globally $\alpha$-H\"older continuous.


{{\vskip 2mm \noindent }} (ii) \ for any  $t \ge 0$,
 $\Psi_{\lambda} : \mathbb{R} ^d \to \mathbb{R} ^d$ is a diffeomorphism of class $C^2$.

{{\vskip 2mm \noindent }} (iii)  $\Psi_{\lambda}^{-1}$ has bounded
first and second spatial
 derivatives, uniformly in $t \in [0,+\infty)$.
\end{lemma}

\begin{proof} Assertion (i) follows by Lemma \ref{schaud}.

\vskip 2mm (ii) Recall the classical Hadamard theorem (see, for
instance, \cite[page 330]{protter}): Let $g : \mathbb{R} ^d \to
\mathbb{R} ^d$  be of class $C^2.$ Suppose that $\lim_{|x| \to
\infty} |g(x) | = + \infty$ and that the Jacobian matrix $Dg(x)$
is an isomorphism of $\mathbb{R} ^d$ for all $x \in \mathbb{R}
^d$. Then $g$ is a
 $C^2$-diffeomorphism of $\mathbb{R} ^d$.
Applying this result to $\Psi_{\lambda}$, we get the assertion.

% ci vuole anche la suriettivita'
%Alternatively,  we can observe that by the inverse function
%theorem we need only to prove that, for fixed $t$,
%$\Psi_\lambda(t,\cdot)$ is injective.
% But assume that $x,y \in \mathbb{R}^d$ satisfy $\Psi_\lambda(t,x)=\Psi_\lambda(t,y)$, then
%$
%x-y = \psi_\lambda(t,x)-\psi_\lambda(t,y)
%$ and since $\psi_\lambda(t,\cdot)$ has a bounded gradient and $\|\nabla \psi_\lambda(t,\cdot)\|_0 < 1$ we obtain
%$
%|x-y| \le \|\nabla \psi_\lambda(t,\cdot)\| |x-y|
% $ and thus $x=y$.


\vskip 2mm (iii) We know that $ \Psi_{\lambda}^{-1} $ is of class
$C^2$. Moreover
$$
D \Psi_{\lambda}^{-1}(t,y) = [D\Psi_{\lambda}(t,
\Psi_{\lambda}^{-1} y)]^{-1} = [I + D\psi_{\lambda}(t,
\Psi_{\lambda}^{-1} y)]^{-1}
$$
$$
= \sum_{k \ge 0} (- D\psi_{\lambda}(t, \Psi_{\lambda}^{-1}
y))^k,\;\;\; y \in \mathbb{R} ^d.
$$
It follows that
$$ \sup_{t\geq0}\Vert D \Psi_{\lambda}^{-1}(t, \cdot)
  \Vert_{0} \le \sum_{k \ge 0} \big (\sup_{t\geq0}
   \Vert D\psi_{\lambda}(t, \cdot )
     \Vert_{0} \big )^k < \infty.
$$
 This shows the boundedness of the first derivative.
 Arguing in a similar way we get also the
 boundedness of the second derivative since
 $$
 D^2 \Psi_\lambda^{-1}(t,y) =  -  [D\Psi_{\lambda}(t,
\Psi_{\lambda}^{-1} y)]^{-1} D^2  \Psi_{\lambda}(t, \Psi_{\lambda}^{-1}(t,y)) \left\{ [D\Psi_{\lambda}(t,
\Psi_{\lambda}^{-1} y)]^{-1} \right\}^{\otimes 2} .
 $$
\end{proof}

\vskip 2mm
 In the sequel we will use a value of $\lambda$ for which Lemma \ref{diff} holds
 and simply write $\psi$ and $\Psi$ for $\psi_{\lambda}$ and
$\Psi_{\lambda}$.

\vskip 2mm \textbf{Step 2} (conjugated SDE). Define
\[
\widetilde{b}(t,y)=-\lambda\psi(t,\Psi^{-1}(t,y)),\quad\widetilde{\sigma
}(t,y)=D\Psi(t,\Psi^{-1}(t,y))
\]
and consider, for every $s\in\left[  0,T\right]  $ and
$y\in{\mathbb{R}}^{d}$,
the SDE%
\begin{equation}
Y_{t}=y+\int_{s}^{t}\tilde{\sigma}(u,Y_{u})dW_{u}+\int_{s}^{t}\widetilde
{b}(u,Y_{u})du,\qquad t\in\lbrack s,T]. \label{conjugated SDE}
\end{equation}
This equation is equivalent to equation (\ref{SDE}), in the
following sense.\ If $X_{t}$ is a solution to \eqref{SDE}, then
$Y_{t}=\Psi(t,X_{t})$ verifies equation (\ref{conjugated SDE})
with $y=\Psi(s,x)$: it is sufficient to apply It\^{o} formula to
$\Psi(t,X_{t})$ and use equation (\ref{pde}). Viceversa, given a
solution $Y_{t}$ of equation (\ref{conjugated SDE}), let
$X_{t}=\Psi^{-1}(t,Y_{t})$; let us prove that $X_{t}$ is a
solution of \eqref{SDE} with $x=\Psi^{-1}(s,y)$. In the sequel
summation over repeated indices is understood. We set
  $\rho(t,\cdot)=\Psi^{-1}(t,\cdot)$. Using Lemma \ref{diff},
   a direct application of It\^{o} formula gives
\[
dX_{t}^{i}=\nabla_{j}\rho_{t}^{i}(Y_{t})\,\widetilde{b}_{t}^{j}(Y_{t}%
)dt\,+\,\nabla_{j}\rho_{t}^{i}(Y_{t})\widetilde{\sigma}_{t}^{jk}(Y_{t}%
)dW_{t}^{k}
$$ $$+\frac{1}{2}\nabla_{j}\nabla_{k}\rho_{t}^{i}(Y_{t})\widetilde
{\sigma}_{t}^{j\ell}(Y_{t})\widetilde{\sigma}_{t}^{k\ell}(Y_{t})dt+\partial
_{t}\rho_{t}^{i}(Y_{t})dt
\]
and since $(\nabla\rho_{t})^{ij}(y)=\nabla_{j}\rho_{t}^{i}(y)=[((\nabla
\Psi_{t})(\rho_{t}(y)))^{-1}]^{ij}=[\tilde{\sigma}_{t}(y)^{-1}]^{ij},$
\[
\nabla_{k}\nabla_{\ell}\rho_{t}^{i}(y)=-[\tilde{\sigma}_{t}(y)^{-1}%
]^{im}(\nabla_{a}\nabla_{b}\Psi_{t}^{m})(\rho_{t}(y))[\tilde{\sigma}%
_{t}(y)^{-1}]^{ak}[\tilde{\sigma}_{t}(y)^{-1}]^{b\ell}%
\]
and $\partial_{t}\rho_{t}(y)^{i}=[\tilde{\sigma}_{t}(y)^{-1}]^{ij}\partial
_{t}\Psi_{t}(\rho_{t}(y))^{j}$. The diffusion coefficient in the equation for
$dX$ is $\nabla\rho_{t}(y)\widetilde{\sigma}_{t}(y)=I$ while the computation
of the drift gives:
\[%
\begin{split}
&  \nabla_{j}\rho_{t}(y)^{i}\tilde{b}_{t}(y)^{j}+\frac{1}{2}\nabla_{j}%
\nabla_{k}\rho_{t}^{i}(y)\widetilde{\sigma}_{t}^{j\ell}(y)\widetilde{\sigma
}_{t}^{k\ell}(y)+\partial_{t}\rho_{t}^{i}(y)\\
&  \qquad=[\tilde{\sigma}_{t}(y)^{-1}]^{ij}\left[  -\lambda\psi_{t}(\rho
_{t}(y))^{j}-\frac{1}{2}(\Delta\psi_{t}^{j})(\rho_{t}(y))+\partial_{t}\psi
_{t}(\rho_{t}(y))^{j}\right] \\
&  \qquad=[\tilde{\sigma}_{t}(y)^{-1}]^{ij}\left[  b_{t}(\rho_{t}%
(y))^{k}(\nabla_{k}\Psi_{t}^{j})(\rho_{t}(y))\right] \\
&  \qquad=[\tilde{\sigma}_{t}(y)^{-1}]^{ij}\tilde{\sigma}_{t}(\rho
_{t}(y))^{jk}b_{t}(\rho_{t}(y))^{k}=b_{t}(\rho_{t}(y))^{i}%
\end{split}
\]
so we obtain equation~(\ref{SDE}).

\textbf{Step 3} (proof of (i) and (ii)). We have  clearly
$\widetilde{b}$ and $\widetilde{\sigma}\in C_{b,T}^{1+{\alpha}}$.
By classical results (see \cite[Chapter 2]{K}) this implies
existence and uniqueness of a strong solution $Y$ of equation
(\ref{conjugated SDE}) and even the existence of a
$C^{1,\alpha^{\prime}}$ ($\alpha^{\prime}<\alpha$) stochastic flow
of diffeomorphisms $\varphi_{s,t}$ associated to equation
(\ref{conjugated SDE}). By the results of steps 1 and 2, these
facts hold true for equation~(\ref{SDE}) and     $\phi_{s,t} =
\Psi^{-1}_t \circ \varphi_{s,t} \circ \Psi_s$ is the corresponding
flow.





\vskip 2mm \textbf{Step 4} (proof of (iii)). Let $\psi^{n}$ be the
solution of the parabolic problem associated to $b_{n}$. Since
$b^{n}\rightarrow b$ in $C_{b,T}^{\alpha}$. By Lemma~\ref{schaud}
we have $\psi^{n}\rightarrow\psi$ in $C_{b,T}^{2+{\alpha}}$.
Consider the flows $\varphi_{s,t}^{n}=\Psi
_{t}^{n}\circ\phi_{s,t}^{n}\circ(\Psi_{s}^{n})^{-1}$ which satisfy
\begin{equation}
\varphi_{s,t}^{n}(y)=y+\int_{s}^{t}\widetilde{b}_{u}^{n}\circ\varphi
_{s,u}^{n}(y)du+\int_{s}^{t}\widetilde{\sigma}_{u}^{n}\circ\varphi
_{s,u}^{n}(y)\cdot dW_{u},
\end{equation}
We have
$\widetilde{\sigma}_{u}^{n}\rightarrow\widetilde{\sigma}_{u}$
 and $\widetilde{b}_{u}^{n}\rightarrow\widetilde {b}_{u}$ in
$C_{b,T}^{1+{\alpha}}$. By standard argument using the Gronwall
lemma, the Doob inequality and the Burkholder inequality
%and Kolmogorov lemma
%
%
%
%
%
%
%
(compare, for instance, with the proof of \cite[Theorem
II.2.1]{K})) we obtain \eqref{stability1}.
%By an easy comparison
%argument based on eqn.~(\ref{eq:approx-phi})
%and~(\ref{eq:approx-Dphi}) we obtain that $\phi^n\to \phi$ in
%$C^{1}$
%
%
%
%
%
%
%

Concerning the derivative of the stochastic flow, we first prove
that
 there exists $C>0$ such that
\begin{equation}
\sup_{x\in{\mathbb{R}}^{d}} \sup_{0\leq s \leq T}  E[\sup_{u \in [s,T]}\Vert D\phi_{s,u}%
^{n}(x)\Vert^{p}]\leq C, \label{bound}
\end{equation}
for any $n\in \mathbb{N}$ and $p\geq 1$ (here $\Vert\cdot\Vert$
denotes the Hilbert-Schmidt norm).

 Arguing as in the proof of \cite[Theorem II.3.1]{K}, we get the
 following linear equation for the derivative $D\phi_{s,t}(x)$
\begin{equation} \label{dove}
\begin{split}
 \lbrack D\Psi_{t}(\phi_{s,t}(x))] D\phi_{s,t}(x)  &  =D\Psi_{s}(x)
  +\int_{s}
^{t}[D^{2}\Psi_{u} (\phi_{s,u}(x))]D\phi_{s,u}(x)dW_{u}\\
& -\lambda\int_{s}^{t}[D\psi_{u}(\phi_{s,u}(x))]D\phi_{s,u}(x)du,
\end{split}
\end{equation}
$ 0 \le s\le t \le T, \; x \in \mathbb{R}^d.$ Thanks to
\eqref{stima}, we get that  $\psi^{n}\rightarrow\psi$ in
$C_{b,T}^{2+{\alpha}}$. Using this fact together with
\eqref{bound} and \eqref{dove}, we finally obtain, for any $p\geq
1$,
\begin{equation} \label{stima1}
\lim_{n\rightarrow\infty}\sup_{x\in{\mathbb{R}}^{d}}\sup_{0\leq
s\leq T} E[\sup_{u \in [s,T]}\Vert
D\phi_{s,u}^{n}(x)-D\phi_{s,u}(x)\Vert^{p}]=0.
\end{equation}
\end{proof}

\vskip 2mm We show now that the inverse flow ${\phi}_{s,t}^{-1}$
is directly related to the solutions of a simple backward
stochastic differential equations, of the same form as the
original one (only the drift has opposite sign). This result will
be useful   to prove uniqueness for our SPDE.


\begin{lemma}\label{fra}
\bigskip
Under  Hypothesis \ref{hy1} the process $\left(
\phi_{s,t}^{-1}\left(  y\right)  \right) _{s\in\left[ 0,t\right]
}$ is the unique solution of the backward SDE
\[
\phi_{s,t}^{-1}\left(  y\right)  =y-\int_{s}^{t}b_{r}\left(  \phi_{r,t}%
^{-1}\left(  y\right)  \right)  dr-\left[  W_{t}\left(
\omega\right) -W_{s}\left(  \omega\right)  \right].
\]
and
\begin{equation}
\sup_{x\in{\mathbb{R}}^{d}} \sup_{0\leq u \leq T}  E[\sup_{s \in [0,u]}\Vert D\phi^{-1}_{s,u}%
(x)\Vert^{p}]<\infty
 \label{bound-inv}
\end{equation}
for any $p\ge 1$.
\end{lemma}

\begin{proof} We have%
\begin{align*}
\phi_{s,t,\omega}\left(  x\right)   &  =x+\int_{s}^{t}b_{r}\left(
\phi_{s,r,\omega}\left(  x\right)  \right)  dr+W_{t}\left(
\omega\right)
-W_{s}\left(  \omega\right) \end{align*}
and then
\begin{align*}
\phi_{s,t,\omega}\left(  \phi_{s,t,\omega}^{-1}\left(  y\right)
\right)   & =\phi_{s,t,\omega}^{-1}\left(  y\right)
+\int_{s}^{t}b_{r}\left( \phi_{s,r,\omega}\left(
\phi_{s,t,\omega}^{-1}\left(  y\right)  \right) \right)
dr+W_{t}\left(  \omega\right)  -W_{s}\left(  \omega\right)  .
\end{align*}
But
$
\phi_{s,r,\omega}\circ\phi_{s,t,\omega}^{-1}=\phi_{r,t,\omega}^{-1}%
$
and thus
\[
y=\phi_{s,t,\omega}^{-1}\left(  y\right)  +\int_{s}^{t}b_{r}\left(
\phi_{r,t,\omega}^{-1}\left(  y\right)  \right)  dr+W_{t}\left(
\omega\right)  -W_{s}\left(  \omega\right)  .
\]
The proof of the bound~(\ref{bound-inv})  is then similar to that
of eq.~(\ref{stima1}) once taken into account the backward
character of the equation.
\end{proof}







\begin{remark}
\label{civuole}  Note that, for any $M > 0$ and $s \in[0,T]$, using the
stability estimates in Th.~\ref{th:flow1}
%and Lemma \ref{fra}
and the bounds
\[
E \Big [  \int_{|y|\leq M} \sup_{u\in[s,T]}
 |\phi_{s,u}^{n}(y)-\phi_{s,u}(y)|^{p}dy \, +
\int_{|y|\leq M} \sup_{u\in[s,T]} \Vert D\phi_{s,u} ^{n}(y)- D
\phi_{s,u} (y)\Vert^{p}dy \Big]
\]
$$ \leq C \Big(
\sup_{x\in{\mathbb{R}}^{d}}E[ \sup_{u\in[s,T]}  |\phi_{s,u}^{n}%
(x)-\phi_{s,u}(x)|^{p}] +  \sup_{x\in{\mathbb{R}}^{d}}
 E[\sup_{u \in [s,T]}\Vert D\phi_{s,u}
^{n}(x)- D \phi_{s,u} (x)\Vert^{p}] \Big),
$$
 with $p \ge 1$, and possibly passing
 to a subsequence, still denoted by $\phi^{n}$, we have that
$\phi_{s,t}^{n}(x)$ and $D\phi_{s,t}^{n}(x)$ converge
 respectively
  to $\phi_{s,t}(x)$ and $D\phi_{s,t}(x)$, a.s in $x$ with
$|x|\leq M$,  $P$-a.s. in $\omega\in\Omega$ and uniformly  in
$t\in\lbrack s,T]$.
\end{remark}

In the sequel, we will often use the standard \textit{mollifiers
approximation.} Let $\vartheta: {\mathbb{R}} ^{d} \to{\mathbb{R}}
$ be a test function such that $0 \le\vartheta(x) \le1$, $x
\in{\mathbb{R}} ^{d}$, $\vartheta (x)= \vartheta (-x)$,
$\int_{{\mathbb{R}} ^{d}} \theta(x)dx =1$, Supp$(\vartheta)
\subset$ $B(2)$, $\vartheta(x) =1, $ $x \in B(1)$ ($B(r)$ denotes
the euclidean ball centered in 0 and with radius $r>0$). For any
${\varepsilon} >0$, define $(\vartheta_{{\varepsilon} }) $,
$\vartheta_{{\varepsilon} } (z) = {{\varepsilon} }^{-d}
\vartheta(z / {\varepsilon} )$. Let $g : {\mathbb{R}} ^{d}
\to{\mathbb{R}} ^{n}$ be a measurable function. The mollifiers
approximation of $g$ is
\begin{equation} \label{molli}
g^{{\varepsilon} }(x) = g * \vartheta_{{\varepsilon} } (x) =
\int_{{\mathbb{R}} ^{d}} g(x-y) \vartheta_{{\varepsilon} } (y)
dy,\;\;\; x \in{\mathbb{R}} ^{d}.
\end{equation}
If $g$ depends also on time $t$, we consider $g^{\varepsilon}(t,
x) = (g(t, \cdot) * \vartheta_{{\varepsilon} }) (x)$, $t
\in[0,T]$, $x \in \mathbb{R} ^d$.

{{\vskip2mm \noindent}} Let $\phi_{t} = \phi_{0,t}$ be the
stochastic flow associated to equation (\ref{SDE}) starting at
$0$. Let us  define  the Jacobian determinant of the flow $\phi $:
\[
J_{s,t}(x) = \det [D \phi_{s,t}(x)],\;\; 0 \le s \le t \le T, \;\;
\mbox{and}\;\; J_{t}(x) = \det [D \phi_{t}(x)],\;\; t \in [0,T].
\]



We finish the section with a result of independent
interest. It concerns a special situation when the given vector field $b$ is
of zero distributional divergence and give rise to a measure-preserving flow.

\begin{lemma}
\label{lemma change} Assume ${\mathrm{div}\,}b(t,\cdot)=0$,
$t\in\lbrack0,T],$ in the sense of distributions. Then the
stochastic flow $\phi$ is $P$-a.s. measure-preserving, i.e.,
$J_{s,t}(x)=1$, for all $0\leq s\leq t\leq T$ and all
$x\in{\mathbb{R}}^{d}$, $P$-a.s..
\end{lemma}

\begin{proof}
Let $b_{n}$ be $C_{b,T}^{\infty}$ vector fields that converges to
$b$ in $C_{b,T}^{{\alpha}}$ and such that
${\mathrm{div}\,}b_{n}=0$. The functions $b_{n}$ can be
constructed as in \eqref{molli}.  Let $\phi^{n}$ be the associated
smooth diffeomorphism. Applying \cite[Theorem II.3.1]{K} and the
well known Liouville theorem, we get that the diffeomorphism
$\phi^{n}$ preserves the Lebsegue measure since $b_{n}$ is a
divergenceless vector field.
%and the vector fields associated to the noise are
%constant (see [Kunita]).
%
%
%
%
%
%
%
Then $\det [D\phi_{s,t}^{n}(x)]=1$ for all $0\leq s<t\leq T$ and
all $x$, $P$-a.s.. Fix $x\in{\mathbb{R}}^{d}$ and $s \in[0,T]$. By
\eqref{stima1}, there exists a subsequence (possibly depending on
$x,s$ and still denoted by $D\phi_{s,u}^{n}$) such that $P$-a.s.
\[
\sup_{s\leq u\leq T}\Vert D\phi_{s,u}^{n}(x)-D\phi_{s,u}(x)\Vert
^{2}\rightarrow0,\;\;\mbox{as}\;\;n\rightarrow\infty.
\]
We find that $J_{s,t}(x)=1$, for any $x\in{\mathbb{R}}^{d}$,  and
so $\phi_{s,t}$ is a measure-preserving diffeomorphism $P$-a.s.
for any $s<t\leq T$.
\end{proof}

%%%%%%%%%%%%%%%%%%%%%%%%%
%
%
%
%
%
%
%

%%%%%%%%%%%%%%%%%%%%%%%%%%%%%%%%
%
%
%
%
%
%
%

\section{Stochastic transport equation. Existence of weak solutions}

\label{sec:transp}
%Throughout this
%section we assume the condition $b\in C^{{\alpha}}$ for some ${\alpha}>0$ and
%$
%{\mathrm{div}\,}  b \in L^q
%$
% for some $q \ge 1$. We let $p$ the conjugate exponent $1/p = 1-1/q$.
%
%
%
%
%
%
%

%\textbf{METTERE UN REMARK SUL CASO $u_{0} \in L^{p}$}

\begin{definition}
\label{def solution}Given $u_{0}\in L^{\infty}({\mathbb{R}}^{d})$, a weak
$L^{\infty}$-solution of the Cauchy problem (\ref{SPDE})\ is a measurable
function $(u(t,x,\omega),t\geq0,x\in{\mathbb{R}}^{d},\omega\in\Omega)$ such that

\begin{itemize}
\item [(i)]$\sup_{t\in\lbrack0,T],x\in{\mathbb{R}}^{d},\omega\in\Omega}\left|
u(t,x,\omega)\right|  <\infty$

\item[(ii)] for any test function $\theta\in C_{0}^{0}({}{\mathbb{R}}^{d})$,
the process $t\mapsto\int_{{}{\mathbb{R}}^{d}}u(t,x)\theta(x)dx$ is continuous
and ${}\mathcal{F}_{t}$-adapted;

\item[(iii)] for any test function $\theta\in C_{0}^{\infty}({}{\mathbb{R}%
}^{d})$, the process $t\mapsto\int_{{}{\mathbb{R}}^{d}}u(t,x)\theta(x)dx$ is
an ${}\mathcal{F}_{t}$-semimartingale and
\begin{align*}
\int_{{\mathbb{R}}{}^{d}}u(t,x)\theta(x)dx  &  =\int_{{\mathbb{R}}{}^{d}}%
u_{0}(x)\theta(x)dx+\int_{0}^{t}ds\int_{{}{\mathbb{R}}^{d}}u(s,x)[b_{s}%
(x)\cdot\nabla\theta(x) \\ & +{\mathrm{div}\,}b_s(x)\theta(x)]dx
  +\sum_{i=1}^{d}\int_{0}^{t}\left(  \int_{{}{\mathbb{R}}^{d}}u(s,x)D_{i}%
\theta(x)dx\right)  \circ dW_{s}^{i}%
\end{align*}
\end{itemize}
\end{definition}


%The following result is well known.({\bf toglierlo??})
%\begin{proposition}
%\label{lemma:diff-lp} If $\Phi$ is a diffeomorphism and $f\in
%L^{p}(\mathbb{R} ^d)$ then $f\circ\Phi\in L^{p}(\mathbb{R} ^d)$,
%\[
%\int_{\mathbb{R} ^d}|f(\Phi(x))|^{p}dx=\int_{\mathbb{R} ^d}|f(y)|^{p}|J\Phi^{-1}(y)|dy,
%\]%
%\[
%\int_{\mathbb{R} ^d} f(\Phi(x))g(x)dx=\int_{\mathbb{R} ^d}
%f(y)g(\Phi^{-1}(y))|J\Phi^{-1}(y)|dy,
%\]
%for any continuous function $g$ with compact support. Here
%$J\Phi^{-1}$ denotes the Jacobian determinant of $D\Phi^{-1}$
%\end{proposition}

\begin{theorem}
\label{them SPDE n1} Assume  Hypothesis \ref{hy1} and $
\mathrm{div}\,  b \in L_{}^{1}\left(   0,T;
L^1_{loc}(\mathbb{R}^d)\right). $ Given $u_{0}\in
L^{\infty}({}{\mathbb{R}}^{d})$, the stochastic process $u(t,x)$
defined as $u(t,x):=u_{0}(\phi_{t}^{-1}(x))$ is a solution of
(\ref{SPDE}).
\end{theorem}

\begin{proof}
The function $u_{0}(\phi_{t}^{-1}(x))$ is bounded and measurable
in $\left( t,x,\omega\right)  $. If $\theta \in C^0_0(\mathbb{R}
^d)$, by changing variable we have:
%from
% Proposition~\ref{lemma:diff-lp},
\[
\int_{{}{\mathbb{R}}^{d}}u(t,x)\theta(x)dx=\int_{{}{\mathbb{R}}^{d}}%
u_{0}(y)\theta(\phi_{t}(y))J_{t}(y)dy .
\]
 This easily yields
property (ii) of Definition \ref{def solution}.

To prove property (iii) consider the flow  $\phi
_{t}^{\varepsilon}$ for the regularized vector field
$b^{\varepsilon}$ (see \eqref{molli}) and let
$J_{t}^{\varepsilon}(y)$ be its Jacobian determinant. Note that
$u_0 \circ (\phi^\varepsilon_t)^{-1} \to u_0 \circ \phi_t^{-1}$
weakly in $L^{\infty}(\mathbb{R}^d)$, uniformly in $t \in [0,T]$
and $P$-a.s., indeed for $\theta \in C_0^\infty(\mathbb{R} ^d)$ we
have
$$
\int_{{}{\mathbb{R}}^{d}} (u_0 \circ (\phi^\varepsilon_t)^{-1}
)(y) \theta(y) dy = \int_{{}{\mathbb{R}}^{d}} u_0(y)
\theta(\phi^{\varepsilon}_t(y)) J_t^\varepsilon(y) dy \to
\int_{{}{\mathbb{R}}^{d}} u_0(y) \theta(\phi_t(y)) J_t(y) dy
$$
using Remark~\ref{civuole}. By density we can extend this
convergence to any $\theta \in L^1(\mathbb{R} ^d)$. Moreover since
$b^{\varepsilon}$ is smooth, it is easy to prove that
\[
dJ_{t}^{\varepsilon}(y)={\mathrm{div}\,}b^{\varepsilon}_t(\phi_{t}^{\varepsilon
}(y))J_{t}^{\varepsilon}(y)dt.
\]
and by the It\^o formula we find
\begin{equation} \label{ii}
\begin{split}
\int_{{\mathbb{R}}^{d}}u_{0}(y)\theta(\phi_{t}^\varepsilon
(y))J_{t}^\varepsilon (y)dy  &
=\int_{{\mathbb{R}}^{d}}u_{0}(y)\theta(y)dy+\int_{0}^{t}ds\int_{{\mathbb{R}%
}^{d}}u_{0}(y)L^{b^\varepsilon}\theta(\phi_{s}^\varepsilon (y))J_{s}^\varepsilon(y)dy\\
&
+\int_{0}^{t}ds\int_{{\mathbb{R}}^{d}}u_{0}(y)\theta(\phi_{s}^\varepsilon
(y)){\mathrm{div}\, }b^\varepsilon_s
(\phi_{s}^\varepsilon(y))J_{s}^\varepsilon(y)dy\\
&
+\int_{0}^{t}dW_{s}\cdot\int_{{\mathbb{R}}^{d}}u_{0}(y)\nabla\theta
(\phi_{s}^\varepsilon(y))J_{s}^\varepsilon(y)dy.
\end{split}
\end{equation}

To pass to the limit  as $\varepsilon \to 0$ the only difficulty
is given by the term
$$
\int_{0}^{t}ds\int_{{\mathbb{R}}^{d}}u_{0}(y)\theta(\phi_{s}^\varepsilon
(y)){\mathrm{div}\, }b^\varepsilon_s
(\phi_{s}^\varepsilon(y))J_{s}^\varepsilon(y)dy =
\int_{0}^{t}ds\int_{{\mathbb{R}}^{d}}u_{0}((\phi_{s}^\varepsilon)^{-1}
(y))\theta(y){\mathrm{div}\, }b^\varepsilon_s (y)dy
$$
Now ${\mathrm{div}\, }b^\varepsilon$ converges in  $
L_{}^{1}\left( 0,T; L^1_{loc}(\mathbb{R}^d)\right)$ to
${\mathrm{div}\, }b$
  and, using also the previous weak convergence of
 $u_0 \circ (\phi^\varepsilon_t)^{-1}$,
  we can conclude the convergence, for any fixed $\omega$,
   of this term toward
$$
\int_{0}^{t}ds\int_{{\mathbb{R}}^{d}}u_{0}((\phi_{s})^{-1}
(y))\theta(y){\mathrm{div}\, }b_s (y)dy
$$
as required. For the stochastic integral we have
 $$
\int_{0}^{t}dW_{s}\cdot\int_{{\mathbb{R}}^{d}}u_{0}(y)\nabla\theta
(\phi_{s}^\varepsilon(y))J_{s}^\varepsilon(y)dy \, \to \,
 \int_{0}^{t}dW_{s}\cdot\int_{{\mathbb{R}}^{d}}u_{0}(y)\nabla\theta
(\phi_{s}(y))J_{s}(y)dy
 $$
uniformly on $[0,T]$ in $L^2 (\Omega)$ as $\varepsilon \to 0$.
Then we get
\[%
\begin{split}
\int_{{\mathbb{R}}^{d}}u_{0}(y)\theta(\phi_{t}(y))J_{t}(y)dy  &
=\int_{{\mathbb{R}}^{d}}u_{0}(y)\theta(y)dy+\int_{0}^{t}ds\int_{{\mathbb{R}%
}^{d}}u_{0}(y)L^{b}\theta(\phi_{s}(y))J_{s}(y)dy\\
&  +\int_{0}^{t}ds\int_{{\mathbb{R}}^{d}}u_{0}(y)\theta(\phi_{s}%
(y)){\mathrm{div}\,}b(\phi_{s}(y))J_{s}(y)dy\\
&  +\int_{0}^{t}dW_{s}\cdot\int_{{\mathbb{R}}^{d}}u_{0}(y)\nabla\theta
(\phi_{s}(y))J_{s}(y)dy.
\end{split}
\]

By writing the above formula with $\theta$ replaced by $\nabla_{i}\eta$, where
$\eta$ is a test function, we discover that $\int_{{\mathbb{R}}^{d}}%
u_{0}(y)\nabla\theta(\phi_{s}(y))J_{s}(y)dy$ is a semimartingale and so it can
be integrated in the Stratonovich sense.

Using the Stratonovich version of It\^{o} formula
\[
\theta(\phi_{t}(y))=\theta(y)+\int_{0}^{t}b(\phi_{s}(y))\cdot\nabla\theta
(\phi_{s}(y))ds+\int_{0}^{t}\nabla\theta(\phi_{s}(y))\circ dW_{s}.
\]
and a stochastic Fubini theorem (see \cite{IW}) gives finally
\[%
\begin{split}
\int_{{\mathbb{R}}^{d}}u_{0}(y)\theta(\phi_{t}(y))J_{t}(y)dy  &
=\int_{{\mathbb{R}}^{d}}u_{0}(y)\theta(y)dy\\ & +\int_{0}^{t}ds\int_{{\mathbb{R}%
}^{d}}u_{0}(y)b(\phi_{s}(y))\cdot\nabla\theta(\phi_{s}(y))J_{s}(y)dy\\
&  +\int_{0}^{t}ds\int_{{\mathbb{R}}^{d}}u_{0}(y)\theta(\phi_{s}%
(y)){\mathrm{div}\,}b(\phi_{s}(y))J_{s}(y)dy\\
&  +\int_{0}^{t}dW_{s}\circ\int_{{\mathbb{R}}^{d}}u_{0}(y)\nabla\theta
(\phi_{s}(y))J_{s}(y)dy
\end{split}
\]
which is exactly the equation we want to prove.
\end{proof}

\section{Uniqueness of weak solutions}

In the deterministic case with non-regular vector fields $b$, uniqueness of
weak $L^{\infty}$-solutions is proved by means of the concept of renormalized
solutions, see \cite{DiPernaLions}, \cite{Ambrosio}. The technical tool is a
commutator lemma, where the role of the assumptions on $b$ is immediately visible.

One of the usual forms of the commutator lemma is the following
one, taken from \cite{Lions}, Chapter 2. Let
$\vartheta_{\varepsilon}$ be the mollifiers introduced in
\eqref{molli}; the convolution $\vartheta_{\varepsilon}\ast f$ $=
\int\vartheta _{\varepsilon}\left( \cdot-x\right) f\left( x\right)
dx$ is  extended also to
distributions $f$.

\vskip 1mm Let $\eta : \mathbb{R}^d \to \mathbb{R}^d$ be a
$C^1$-diffeomorphism; in this section, we will  use the
 following notation for the Jacobian determinant
$$
J \eta (x) = \text{det}[D \eta (x) ],\;\;\; x \in \mathbb{R}^d.
$$

Given $g\in L_{loc}^{\infty}\left(  \mathbb{R}%
^{d},\mathbb{R}\right)  $, $v\in L_{loc}^{1}\left(  \mathbb{R}^{d}%
,\mathbb{R}^{d}\right)  $ , $\mathrm{div}\,  v\in
L_{loc}^{1}\left( \mathbb{R}^{d},\mathbb{R}\right)  $, denote by
$v\cdot\nabla g$ the distribution defined on smooth compact
support test functions $f$ as
\[
\left(  v\cdot\nabla g\right)  \left(  f\right)  :=-\int
gv\cdot\nabla fdx-\int gf\mathrm{div}\,  v \, dx.
\]
Define also the commutator $\mathcal{R}_{\varepsilon}\left[  g,v\right]
\left(  x\right)  $ as the (smooth) function%
\[
\mathcal{R}_{\varepsilon}\left[  g,v\right]  \left(  x\right)  =\left[
\vartheta_{\varepsilon}\ast\left(  v\cdot\nabla g\right)  -v\cdot\nabla\left(
\vartheta_{\varepsilon}\ast g\right)  \right]  \left(  x\right)  .
\]

\begin{lemma}
Given
\[
v\in W^{1,1}\left(  \mathbb{R}^{d},\mathbb{R}^{d}\right)  ,\quad g\in
L^{\infty}\left(  \mathbb{R}^{d},\mathbb{R}\right)  ,
\]
we have
\[
\int\left|  \mathcal{R}_{\varepsilon}\left[  g,v\right]  \left(  x\right)
\right|  dx\leq C\left\|  v\right\|  _{W^{1,1}}\left\|  g\right\|
_{L^{\infty}}%
\]
for some $C>0$ independent of $\varepsilon$, $v$ and $g$. Moreover,
\[
\lim_{\varepsilon\rightarrow0}\int\left|  \mathcal{R}_{\varepsilon}\left[
g,v\right]  \left(  x\right)  \right|  dx=0.
\]
\end{lemma}

The strong convergence of commutators requires some weak form of
differentiability of $v$ or $g$. We do not impose it, so we have to use a
relaxed form of this result. First, essentially with the same proof of the
previous lemma one gets the following fact: distributional convergence of
commutators holds under very weak conditions.

{\it Distributional convergence is not sufficient to study the
 deterministic PDE, while we can work with it in the stochastic
case, as explained below.}

% Denote
% by $W_{r}^{1,1}\left(  \mathbb{R}^{d},\mathbb{R}\right)  $ (for $r>0$)
% the space of $W^{1,1}\left(  \mathbb{R}^{d},\mathbb{R}\right)  $-functions

For any $r>0$ denote by $B\left(  r\right)  $ is the
ball centered at the origin of radius $r$ and by $C_{r}^{\infty}(  \mathbb{R}^{d},\mathbb{R})  $ the space
of smooth functions with compact support in $B\left(  r\right)  $.
Moreover, denote by
% \begin{equation} \label{rr}
%\| \cdot \|_{L_{r}^{1}}, \;\;\;  \| \cdot\|_{L^{\infty}_r}
 %\end{equation}
$\| \cdot \|_{L_{r}^{1}}$ and $\| \cdot\|_{L^{\infty}_r}$ respectively,  the $L^{1}$-norm and the $L^{\infty}$-norm on
$B\left( r+1\right) $.


\begin{lemma}
Given
\begin{align*}
v  &  \in L_{loc}^{\infty}\left(  \mathbb{R}^{d},\mathbb{R}^{d}\right)
,\quad{\mathrm{div}}v\in L_{loc}^{1}\left(  \mathbb{R}^{d},\mathbb{R}\right),
\qquad
g    \in L_{loc}^{\infty}\left(  \mathbb{R}^{d},\mathbb{R}\right)
%,\quad \rho\in W_{loc}^{1,1}\left(
%\mathbb{R}^{d},\mathbb{R}\right) \cap L_{loc}^{\infty}\left(
%\mathbb{R}^{d},\mathbb{R}\right)
\end{align*}
for any $\rho \in C_r^\infty(  \mathbb{R}^{d},\mathbb{R})$ and
for sufficiently small $\varepsilon$ we have
\[
\left|  \int\mathcal{R}_{\varepsilon}\left[  g,v\right]  \left(  x\right)
\rho\left(  x\right)  dx\right|  \leq2\left\|  g\right\|  _{L_{r}^{\infty}%
}\left[  \left\|  {\mathrm{div}}v\right\|  _{L_{r}^{1}}\left\|  \rho\right\|
_{L_{r}^{\infty}}+\left\|  v\right\|  _{L_{r}^{\infty}}\left\|  \nabla
\rho\right\|  _{L_{r}^{1}}\right]
\]
and
\[
\lim_{\varepsilon\rightarrow0}\int\mathcal{R}_{\varepsilon}\left[  g,v\right]
\left(  x\right)  \rho\left(  x\right)  dx=0.
\]
\end{lemma}

\begin{proof}
The proof proceeds as the one of the previous lemma from \cite{Lions}, Chapter
2: we prove the inequality for regular fields, then one can extend to
non-regular ones and prove the convergence again first in the regular case. We
omit the details of the second part and only prove the inequality.

\textbf{Step 1}. Let us rewrite the expressions is a suitable way.
\begin{equation*}
\begin{split}
\int & \mathcal{R}_{\varepsilon}\left[  g,v\right] (  x)  \rho(x) dx  = \int \left[
\vartheta_{\varepsilon}\ast\left(  v\cdot\nabla g\right)  -v\cdot\nabla\left(
\vartheta_{\varepsilon}\ast g\right)  \right] (  x) \rho(x) dx
\\
& = \iint [-g(x') v(x') \nabla_{x'} \vartheta_{\varepsilon}(x-x')
- g(x') \vartheta_{\varepsilon}(x-x') \mathrm{div}\,  v(x') ]
\rho(x) \, dx dx'
\\ & \qquad +
 \iint [g(x') v(x)  \vartheta_{\varepsilon}(x-x') \nabla_{x}\rho(x) + g(x')  \rho(x)  \vartheta_{\varepsilon}(x-x') \mathrm{div}\,  v(x) ] \, dx dx'
\end{split}
\end{equation*}
by using $\nabla_{x'} \vartheta_{\varepsilon}(x-x')  = -\nabla_{x}
\vartheta_{\varepsilon}(x-x') $ and integrating by parts in $x$
the first term we get
\begin{equation*}
\begin{split}
 & = \iint g(x')   \vartheta_{\varepsilon}(x-x')\nabla_{x}\rho(x)   [v(x)-v(x')]   \, dx dx'
\\ & \qquad +
 \iint g(x')  \rho(x)  \vartheta_{\varepsilon}(x-x')[ \mathrm{div}\,  v(x)-\mathrm{div}\,  v(x') ] \, dx dx'
\end{split}
\end{equation*}

\textbf{Step 2}. Assume $\varepsilon$ so small that the support of
$\vartheta_{\varepsilon}$ has diameter less than one. We have (using standard
estimates on convolutions)
\[
\left|\int  \mathcal{R}_{\varepsilon}\left[  g,v\right] (  x)  \rho(x) dx \right|
 \leq  2\left\|  g\right\|  _{L_{r}^{\infty}} \left\|
v\right\|  _{L_{r}^{\infty}}\left\|  \nabla\rho\right\|
_{L_{r}^{1}} + 2 \left\|  \rho\right\| _{L_{r}^{\infty}}\left\|
g\right\|  _{L_{r}^{\infty}}\left\| \mathrm{div}\,  v\right\|
_{L_{r}^{1}}
\]%

The inequality is proved. The second assertion follows using
standard properties of mollifiers.  The proof is complete.
\end{proof}

\vskip 2mm In the proof of the uniqueness theorem the
distributional convergence of commutators is sufficient but the
commutator must be composed with the flow. This strategy would
work in the deterministic case as well. But the main difficulty is
that composition with the flow put into play the first derivatives
of the Jacobian of the flow. This would impose too much regularity
on $b$ in the deterministic case. On the contrary, in the
stochastic case, we have a weakly differentiable Jacobian of the
flow under mild assumptions on $b$, as explained in the previous
section.

First, let us write the previous weak commutator lemma in the case when
commutators are applied to a diffeomorphism.


\begin{corollary}
\label{corollary commutator}Let $\phi$ be a $C^{1}$ diffeomorphism of
$\mathbb{R}^{d}$. Assume
\begin{align*}
v  &  \in L_{loc}^{\infty}\left(
\mathbb{R}^{d},\mathbb{R}^{d}\right) ,\quad{\mathrm{div}\, }v\in
L_{loc}^{1}\left(  \mathbb{R}^{d},\mathbb{R}\right),\;\;\; g \in
L_{loc}^{\infty}\left( \mathbb{R}^{d},\mathbb{R}\right)\\ & \quad
J\phi^{-1}\in W_{loc}^{1,1}\left(\mathbb{R}^{d},\mathbb{R}\right).
\end{align*}
then for any $\rho \in C_r^\infty(  \mathbb{R}^{d},\mathbb{R})$
\[
\lim_{\varepsilon\rightarrow0}\int\mathcal{R}_{\varepsilon}\left[  g,v\right]
\left(  \phi\left(  x\right)  \right)  \rho\left(  x\right)  dx=0.
\]
\end{corollary}

\begin{proof}
%For every $\rho\in C_{r}^{\infty}\left(  \mathbb{R}^{d},\mathbb{R}\right)  $,
%Assume that the support of $\rho$ is contained in $B(r)$ for some $r > 0$ and
Set
$
R_{r,\phi}:=\left|  \phi\left(  0\right)  \right|  +\left(  r+1\right)
\left\|  D\phi\right\|  _{L_{r}^{\infty}}
$.
Using
\[
\left|  \phi\left(  x\right)  \right|  \leq\left|  \phi\left(  0\right)
\right|  +\left\|  D\phi\right\|  _{L_{\rho}^{\infty}}\left|  x\right|
\leq\left|  \phi\left(  0\right)  \right|  +r\left\|  D\phi\right\|
_{L_{\rho}^{\infty}}%
\]
we  check that the map $\phi$ send the support of $\rho$ strictly in the ball of radius
$R_{r,\phi}$.
%
%
%Then there exists  a neighborhood  of the support of the function
%$\rho\circ\phi^{-1}$ that is contained in $B\left(
%R_{r,\phi}\right)  $ and for sufficiently small $\varepsilon$ we
%have
%\begin{align*}
%\left|  \int\mathcal{R}_{\varepsilon}\left[  g,v\right]  \left(
%\phi\left( x\right)  \right)  \rho\left(  x\right)  dx\right|
%\leq2\left\| g\right\|  _{L_{R_{r,\phi}}^{\infty}}\left\|
%{\mathrm{div}\, }v\right\| _{L_{R_{r,\phi}}^{1}}\left\|  \rho\right\|
%_{L_{r}^{\infty}}\left\|
%J\phi^{-1}\right\|  _{L_{R_{r,\phi}}^{\infty}}\\
%  +2\left\|  g\right\|  _{L_{R_{r,\phi}}^{\infty}}\left\|  v\right\|
%_{L_{R_{r,\phi}}^{\infty}}\left( \left\| D\phi^{-1}
%  \right\|_{L_{R_{r,\phi}}^{\infty}}  \left\|  \nabla\rho\right\|_{L_{r}^1}
%  +\left\|  \rho\right\|  _{L_{r}^{\infty}}\left\|  \nabla
%J\phi ^{-1}\right\|  _{L_{R_{r,\phi}}^{1}}\right)
%\end{align*}
%and
Standing this preliminary fact, we simply have
\[
\int\mathcal{R}_{\varepsilon}\left[  g,v\right]  \left(  \phi\left(  x\right)
\right)  \rho\left(  x\right)  dx=\int\mathcal{R}_{\varepsilon}\left[
g,v\right]  \left(  y\right)  \rho_{\phi}\left(  y\right)  dx
\]
where the function
$
\rho_{\phi}\left(  y\right)  :=\rho\left(  \phi^{-1}\left(  y\right)  \right)
J\phi^{-1}\left(  y\right)
$
has the support strictly contained in the ball of radius $R_{r,\phi}$. Thus by
the previous lemma
$$
\left|  \int\mathcal{R}_{\varepsilon}\left[  g,v\right]  \left(
\phi\left( x\right)  \right)  \rho\left(  x\right)  dx\right|
\leq2\left\|  g\right\| _{L_{R_{r,\phi}}^{\infty}}\Big[  \left\|
{\mathrm{div}\, }v\right\|
_{L_{R_{r,\phi}}^{1}}\left\|  \rho_{\phi}\right\|  _{L_{R_{r,\phi}}^{\infty}%
}$$ $$  +\left\|  v\right\|  _{L_{R_{r,\phi}}^{\infty}}\left\|
\nabla\rho_{\phi }\right\|  _{L_{R_{r,\phi}}^{1}}\Big]  .
$$
To conclude, it is sufficient to notice that
$
\left\|  \rho_{\phi}\right\| _{L_{R_{r,\phi}}^{\infty}}\leq
\left\| \rho\right\| _{L_{r}^{\infty}}\left\|
J\phi^{-1}\right\| _{L_{R_{r,\phi }}^{\infty}}
$
and, denoting by $[D\phi^{-1}(y)]^*$ the adjoint matrix of
 $[D\phi^{-1}(y)]$,
\begin{align*}
\left\|  \nabla\rho_{\phi}\right\|  _{L_{R_{r,\phi}}^{1}}
\leq\left\| [D\phi^{-1}(\cdot)]^* \left(
\nabla\rho\circ\phi^{-1}\right)(\cdot) \, J\phi^{-1}\left(
\cdot\right) \right\| _{L_{R_{r,\phi}}^{1}}\\
+\left\| \left(
\rho\circ\phi^{-1}\right) \cdot\nabla J\phi^{-1}\right\|
_{L_{R_{r,\phi}}^{1}}
 \\
 \leq  \, \left\| D\phi^{-1}
  \right\|_{L_{R_{r,\phi}}^{\infty}}   \, \left\|
\nabla\rho\right\| _{L_{r}^{1}}+
   \left\|
\rho\right\|  _{L_{r}^{\infty}}\left\|  \nabla J\phi^{-1}\right\|
_{L_{R_{r,\phi}}^{1}}.
\end{align*}
The proof is complete.
\end{proof}

\vskip 2mm We can now prove the uniqueness of weak
$L^{\infty}$-solutions. We postpone
 the main estimate on the derivatives
 of the Jacobian of the flow to the next section.

\begin{theorem}
\label{thm SPDE n2} Under Hypothesis \ref{hy1} and \ref{hy2}, for every $u_{0}\in
L^{\infty}$ there exists a unique weak $L^{\infty}$-solution of
the Cauchy problem (\ref{SPDE}). Thus every solution
$u(t,x,\omega)$ has the form
$u(t,x,\omega):=u_{0}({\phi}_{t}^{-1}\left(  \omega\right)  x)$.
\end{theorem}

\begin{proof}
\textbf{Step 1}. By linearity we have to prove that a solution with initial
condition $u_{0}=0$ vanishes identically. Let us denote by $u$ such a solution.

For $y\in{\mathbb{R}}^{d}$ fixed, ${\varepsilon}>0$, let us choose
the test function $\theta(x):=\vartheta_{{\varepsilon}}\left(
y-x\right)  $ in property (iii) of Definition \ref{def solution}.
Let us define $u_{{\varepsilon}}(t,\cdot) = \vartheta_{{\varepsilon
}} * u(t,\cdot)$.
 We get
\[
u_{{\varepsilon}}(t,y)=u_{0,{\varepsilon}}(y)+\int_{0}^{t}A_{\varepsilon
}\left(  s,y\right)  ds+\sum_{i=1}^{d}\int_{0}^{t}B_{\varepsilon}^{\left(
i\right)  }\left(  s,y\right)  \circ dW_{s}^{i}%
\]
where
\begin{align*}
A_{\varepsilon}\left(  t,y\right)   &  :=\int_{{}{\mathbb{R}}^{d}%
}u(t,x)\left\{  b_{t}(x)\cdot\nabla_{x}[\vartheta_{{\varepsilon}%
}(y-x)]+{\mathrm{div}\, \,}b_{t}(x)\vartheta_{\varepsilon}(y-x)\right\}  dx\\
B_{\varepsilon}^{\left(  i\right)  }\left(  t,y\right)   &  :=\int
_{{}{\mathbb{R}}^{d}}u(t,x)D_{x_{i}}[\vartheta_{{\varepsilon}}\left(
y-x\right)  ]dx.
\end{align*}
All these functions of $y$, namely $u_{{\varepsilon}}(t,y)$,
$u_{0,{\varepsilon}}(y)$, ... , $B_{\varepsilon}^{\left(  i\right)  }\left(
t,y\right)  $, are continuous in $t$, adapted, smooth (of class $C^{3}$ is
required for the next computation) in $y$ (from (i) and (ii) of definition
\ref{def solution}). By the  Stratonovich version of Kunita-It\^{o}-Wentzel formula
(see~\cite[Th. 8.3 page 188]{K})
\begin{align*}
du_{{\varepsilon}}(t,{\phi}_{t}(x))  &  =A_{\varepsilon}\left(  t,{\phi}%
_{t}(x)\right)  dt+\sum_{i=1}^{d}B_{\varepsilon}^{\left(  i\right)  }\left(
t,{\phi}_{t}(x)\right)  \circ dW_{t}^{i}\\
&  +\left(  b\cdot\nabla u_{{\varepsilon}}\right)  (t,{\phi}_{t}%
(x))dt+\sum_{i=1}^{d}\left(  D_{i}u_{{\varepsilon}}\right)  (t,{\phi}%
_{t}(x))\circ dW_{t}^{i}%
\end{align*}
But
\[
\left(  D_{i}u_{{\varepsilon}}\right)  (t,{y})={\int_{{\mathbb{R}}^{d}}%
}u(t,x)D_{y_{i}}[\vartheta_{{\varepsilon}}\left(  y-x\right)  ]\,dx=-{\int
_{{\mathbb{R}}^{d}}}u(t,x)D_{x_{i}}[\vartheta_{{\varepsilon}}\left(
y-x\right)  ]\,dx
\]
hence
\[
du_{{\varepsilon}}(t,{\phi}_{t}(x))=\left[  A_{\varepsilon}\left(  t,{\phi
}_{t}(x)\right)  +\left(  b\cdot\nabla u_{{\varepsilon}}\right)  (t,{\phi}%
_{t}(x))\right]  dt
\]
namely (recall that the initial condition is zero)
\[
u_{{\varepsilon}}(t,{\phi}_{t}(x))=-
 \int_{0}^{t}\mathcal{R}_{\varepsilon }\left[
u_{s},b_{s}\right] \left(  {\phi}_{s}(x)\right)  ds
\]
where $\mathcal{R}_{\varepsilon}\left[  u_{s},b_{s}\right]  \left(
{y}\right)  $ is the commutator defined above.

\vskip 2mm \textbf{Step 2}. For every $\rho\in
C_{r}^{\infty}({}{\mathbb{R}}^{d})$ (see the definition above) we
have
\[
\int u_{{\varepsilon}}(t,{\phi}_{t}(x))\rho\left(  x\right)  dx=\int
u_{{\varepsilon}}(t,{y})\rho\left(  {\phi}_{t}^{-1}(y)\right)  J{\phi}%
_{t}^{-1}(y)dy.
\]
Given $t\in\left[  0,T\right]  $, with probability one, $u_{{\varepsilon}%
}(t,{\cdot})$ converges weak-$*$  to $u(t,{\cdot})$ as $\varepsilon
\rightarrow0$. With probability one, the function $y\mapsto\rho(  {\phi
}_{t}^{-1}(y))  $ is at least continuous with compact support, and also
the function $y\mapsto J{\phi}_{t}^{-1}(y)$ is continuous. Hence, $P$-a.s.,
\begin{align*}
\lim_{\varepsilon\rightarrow0}\int u_{{\varepsilon}}(t,{\phi}_{t}%
(x))\rho\left(  x\right)  dx  &  =\int u(t,{y})\rho\left(  {\phi}_{t}%
^{-1}(y)\right)  J{\phi}_{t}^{-1}(y)dy\\
&  =\int u(t,{\phi}_{t}(x))\rho\left(  x\right)  dx.
\end{align*}
Therefore $\int\left(  \int_{0}^{t}\mathcal{R}_{\varepsilon}\left[
u_{s},b_{s}\right]  \left(  {\phi}_{s}(x)\right)  ds\right)  \rho\left(
x\right)  dx$ converges $P$-a.s. as $\varepsilon\rightarrow0$ and we have
\[
\int u(t,{\phi}_{t}(x))\rho\left(  x\right)  dx=\lim_{\varepsilon\rightarrow
0}\int\left(  \int_{0}^{t}\mathcal{R}_{\varepsilon}\left[  u_{s},b_{s}\right]
\left(  {\phi}_{s}(x)\right)  ds\right)  \rho\left(  x\right)  dx.
\]
The function $\left(  s,x\right)  \mapsto\mathcal{R}_{\varepsilon}\left[
u_{s},b_{s}\right]  \left(  {\phi}_{s}(x)\right)  \rho\left(  x\right)  $ is
continuous with compact support, $P$-a.s., hence
\[
\int u(t,{\phi}_{t}(x))\rho\left(  x\right)  dx=\lim_{\varepsilon\rightarrow
0}\int_{0}^{t}\left(  \int\mathcal{R}_{\varepsilon}\left[  u_{s},b_{s}\right]
\left(  {\phi}_{s}(x)\right)  \rho\left(  x\right)  dx\right)  ds.
\]
{\it If we prove that, given $\rho\in
C_{0}^{\infty}({}{\mathbb{R}}^{d})$ and $t\in\left[  0,T\right] $,
this $P$-a.s.-limit (which exists) is zero, then we have $u$
identically zero (because ${\phi}_{t}$ is a bijection).
}

\vskip 2mm \textbf{Step 3}. Let us check, by means of Corollary
\ref{corollary commutator}, that $$ s\mapsto\int_{\mathbb{R} ^d}
\mathcal{R}_{\varepsilon}\left[  u_{s} ,b_{s}\right] \left(
{\phi}_{s}(x)\right)  \rho\left(  x\right) dx
$$ satisfies
the assumptions of Lebesgue dominated convergence theorem
($P$-a.s.) on $[0,T]$.
 % the convergence to zero is similar.
 This will complete
the proof. For every $s\in\left[  0,T\right]  $ we have (recall
the notation $R_{r,{\phi }}$ from Corollary \ref{corollary
commutator})



\begin{align*}
\left|  \int\mathcal{R}_{\varepsilon}\left[  u_{s},b_{s}\right]
\left( {\phi}_{s}\left(  x\right)  \right)  \rho\left(  x\right)
dx\right| \leq2\left\|  u_{s}\right\|
_{L_{R_{r,{\phi}_{s}}}^{\infty}}\left\| {\mathrm{div}\,
}b_{s}\right\|  _{L_{R_{r,{\phi}_{s}}}^{1}}\left\| \rho\right\|
_{L_{r}^{\infty}}\left\|  J{\phi}_{s}^{-1}\right\|  _{L_{R_{r,{\phi}_{s}}%
}^{\infty}}\\
  +2\left\|  u_{s}\right\|  _{L_{R_{r,{\phi}_{s}}}^{\infty}}\left\|
b_{s}\right\|  _{L_{R_{r,{\phi}_{s}}}^{\infty}}\Big( \left\|
D\phi^{-1}_s
  \right\|_{L_{R_{r,\phi_s}}^{\infty}} \left\|
\nabla \rho\right\|  _{L_{r}^{1}}+\left\| \rho\right\| _{L_{r
}^{\infty}}\left\|  \nabla J{\phi}_{s}^{-1}\right\|  _{L_{R_{r,{\phi}_{s}}%
}^{1}}\Big)
\end{align*}
\[
\leq C_{u,\rho}\Big(  1+\left\|  {\mathrm{div}\, }b_{s}\right\|
_{L_{R_{r,{\phi }_{s}}}^{1}}\left\|  J{\phi}_{s}^{-1}\right\|
_{L_{R_{r,{\phi}_{s}}}^{\infty }}+ \left\| D\phi^{-1}_s
  \right\|_{L_{R_{r,\phi_s}}^{\infty}}
 + \left\|  \nabla J{\phi}_{s}^{-1}\right\|  _{L_{R_{r,{\phi}_{s}}}^{1}%
}\Big)
\]
because of the global boundedness of $u$, $b$, $\rho$.

We know, from the properties of the stochatic flow ${\phi}$, that,
for any $R>0$,
\[
P\left(  \sup_{s\in\left[  0,T\right]  }\left(  \left\|
D{\phi}_{s}\right\| _{L_{R}^{\infty}}+\left\|  {\phi}_{s}\right\|
_{L_{R}^{\infty}}\right) <\infty\right)  =1
\]
hence (recall the definition of $R_{r,{\phi}_{s}}$) there exists a positive
random variable $X$, $P$-a.s. finite, such that
\[
\left\|  {\mathrm{div}\, }b_{s}\right\|
_{L_{R_{r,{\phi}_{s}}}^{1}}\left\| J{\phi}_{s}^{-1}\right\|
_{L_{R_{r,{\phi}_{s}}}^{\infty}}+ \left\| D\phi^{-1}_s
  \right\|_{L_{R_{r,\phi_s}}^{\infty}}
 \] \[ +   \left\|  \nabla
J{\phi}_{s}^{-1}\right\|  _{L_{R_{r,{\phi}_{s}}}^{1}}\leq\left\|
{\mathrm{div}\, }b_{s}\right\|  _{L_{X}^{1}}\left\|
J{\phi}_{s}^{-1}\right\| _{L_{X}^{\infty}}+ \left\| D\phi^{-1}_s
  \right\|_{L_{X}^{\infty}}+
\left\|  \nabla J{\phi}_{s}^{-1}\right\|  _{L_{X}^{1}},
\]
uniformly in $s\in\left[  0,T\right]  $, with probability one. The
term $\int_{0}^{T}\left\|  {\mathrm{div}\, }b_{s}\right\|
_{L_{X}^{1}}ds$ is finite ($P$-a.s.). Since $\left(  s,x\right)
\mapsto D{\phi}_{s}^{-1}\left( x\right)  $ is $P$-a.s. continuous,
we also have
$$ \sup_{s\in\left[ 0,T\right]  } \Big( \left\|
J{\phi}_{s}^{-1}\right\|  _{L_{X}^{\infty}} +
 \left\| D\phi^{-1}_s
  \right\|_{L_{X}^{\infty}} \Big) <\infty,
\;\; \text{P-a.s..}
$$
Finally, we just need%
\begin{equation}
\int_{0}^{T}\left\|  \nabla J{\phi}_{s}^{-1}\right\|  _{L_{X}^{1}}%
ds<\infty,\quad P\text{-a.s.} \label{bound on DJ}%
\end{equation}
This property is crucial and will be proved in the next section.
The proof is complete.
\end{proof}

\section{Estimates on derivatives of the Jacobian}

The aim of this section is to prove the estimate (\ref{bound on DJ}).


Since by Lemma \ref{fra} the equation for ${\phi}_{s}^{-1}$ is
equal to the equation for ${\phi }_{s}$ (up to a sign and
inversion of time), for simplicity of notations we prove the
estimate
\begin{equation}
\int_{0}^{T}\left\|  \nabla J{\phi}_{s}\right\|  _{L_{X}^{1}}ds<\infty,\quad
P\text{-a.s.} \label{bound on DJ direct}%
\end{equation}
in place of estimate (\ref{bound on DJ}). Here $X$ is a generic
$P$-a.s. finite random variable. Recall that
$
\| \cdot \|_{L_{X}^{1}}
$
denotes the $L^{1}$-norm  on the ball $B\left( X+1\right) $.


We know that
\[
J{\phi}_{t}\left(  x\right)  =e^{\int_{0}^{t}{\mathrm{div}\,
}b_{s}\left(  {\phi
}_{s}\left(  x\right)  \right)  ds}%
\]
hence%
\[
\nabla J{\phi}_{t}\left(  x\right)  =J{\phi}_{t}\left(  x\right)
\cdot \nabla\int_{0}^{t}{\mathrm{div}\, }b_{s}\left(
{\phi}_{s}\left(  x\right) \right)  ds
\]
if the integral, as a function of $x$, is differentiable (at least
weakly). The first factor, $J{\phi}_{t}\left(  x\right)  $ is
$P$-a.s. bounded in $\left(  t,x\right)  $ on $\left[  0,T\right]
\times B\left(  X +1 \right)  $  (because it is continuous).
Hence, to prove (\ref{bound on DJ}), it is
sufficient to prove%
\begin{equation}
\int_{0}^{T}\left\|  \nabla\int_{0}^{t}{\mathrm{div}\, }b_{s}\left(  {\phi}%
_{s}\left(  \cdot \right)  \right)  ds\right\|
_{L_{X}^{1}}dt<\infty,\quad
P\text{-a.s.} \label{estimate grad div b}%
\end{equation}

Again the basic tool is a Tanaka type formula, based on a result
for parabolic equations proved in \cite[Theorem 1.2]{Kr1} (see
also \cite[Theorem 10.3]{Kry-Ro}).



\vskip 2mm We introduce a  function space  from \cite{Kr1}
 and \cite{Kry-Ro}. For
$p, q \in (1, + \infty)$,
%we set $W^{2}_p = W^2_p (\mathbb{R} ^d).$
% $ $ (1 -
%\triangle )^{-1}$ $(L^p(\mathbb{R} ^d))$. Note that
% $W^{2,p}(\mathbb{R} ^d) \subset H^{2}_p $ by the Calderon-Zygmund theorem
% and that in the Banach space
%  $H^{2}_p$ an equivalent norm is given by $\| D^2 u \|_{L^p}
% $ $+ \| u\|_{L^p}$, for $u \in H^2_p$.
 we consider  the space $H^{2,q}_p(T)$ consisting of all
 functions
 $u$ defined on $[0,T]$ with values in distributions on $\mathbb{R} ^d$
 such that
 $u\in L^q ([0,T]; W^2_p(\mathbb{R} ^d))$ and there exists  $D_t u
  \in L^q ([0,T]; L^p(\mathbb{R} ^d))$.


\begin{theorem} \label{KR}
 Under  Hypothesis \ref{hy1} and \ref{hy2},  in the space $
 H^{2,q}_p(T)$ there exists a unique
 solution to the
 parabolic equation
$$ \left \{ \begin{aligned}
\frac{\partial F}{\partial t}+L^{b}F={\mathrm{div}\, }b,\;\; \; t
\in [0,T[
\\
F(T,x)=0,\;\;\; x \in \mathbb{R} ^d.
\end{aligned}
\right.
$$
%with the properties%
%\[
%F\left(  0,\cdot\right)  \in W_{loc}^{1,1}\left(  \mathbb{R}^{d}\right)
%,\quad F\in L^{1}\left(  0,T;W_{loc}^{1,1}\left(  \mathbb{R}^{d}\right)
%\right)
%\]%
%\[
%\int_{0}^{T}\int_{K}\left|  D^{2}F\left(  s,{x}\right)  \right|
%^{2+\varepsilon}dxds<\infty
%\]
%for some $\varepsilon>0$ and every compact $K\subset\mathbb{R}^{d}$.
\end{theorem}
 Writing $\frac{\partial F}{\partial t}+ \frac{1}{2} \triangle
  F =  - b \cdot \nabla u + {\mathrm{div}\, }b$,
 this results follows from \cite[Theorem 1.2]{Kr1} (using the fact that
  $b$ is globally bounded). We mention   that an extension  of
   \cite[Theorem 1.2]{Kr1} is given in \cite[Theorem
   10.3]{Kry-Ro}.



%\begin{remark}
%\textbf{LA MIA CONGETTURA} \`{e} che serva ${\mathrm{div}}b\in
%L^{2+\varepsilon^{\prime}}\left(  0,T;L_{loc}^{2+\varepsilon^{\prime}}\right)
%$) per un qualche $\varepsilon^{\prime}>0$.
%\end{remark}

\begin{proposition}
 Under  Hypothesis \ref{hy1} and \ref{hy2}, estimate
(\ref{estimate grad div b}) holds true.
\end{proposition}

\begin{proof} The condition $\frac {d}{p} + \frac{2}{q}
   <1$ implies that if $u \in H^{2,q}_p(T)$ then $u$
   can be represented by a real  function $u(t,x) $ continuous
    on $[0,T]
    \times \mathbb{R} ^d$ such that there exists $\nabla u(t,x)$
 and is continuous on $[0,T]
    \times \mathbb{R} ^d$ (see \cite[Lemma 10.1]{Kry-Ro}).

     The following Ito formula holds
     (it is a special case of  \cite[Theorem 3.7]{Kry-Ro})
\begin{equation} \label{ito}
F\left(  t,{\phi}_{t}\left(  x\right)  \right)  =F\left(
0,x\right) +\int_{0}^{t}\nabla F\left(  s,{\phi}_{s}\left(
x\right)  \right)  \cdot dW_{s}+\int_{0}^{t}\left(  \frac{\partial
F}{\partial s}+L_{b}F\right) \left(  s,{\phi}_{s}\left(  x\right)
\right)  ds.
\end{equation}
%In fact the classical It\^{o} formula requires additional regularity, but with
%some effort of estimates and approximation one can prove this identity anche
%check that all terms are well defined (the estimates are similar to those
%developed below). Thus
Thus
\begin{align*}
\nabla\int_{0}^{t}{\mathrm{div}\, }b_{s}\left(  {\phi}_{s}\left(
x\right) \right)  ds    =\nabla\left[  F\left(  t,{\phi}_{t}\left(
x\right)  \right) -F\left(  0,x\right)  -\int_{0}^{t}\nabla
F\left(  s,{\phi}_{s}\left(
x\right)  \right)  dW_{s}\right] \\
  =\nabla F\left(  t,{\phi}_{t}\left(  x\right)  \right)  \cdot D{\phi}%
_{t}\left(  x\right)  -\nabla F\left(  0,x\right)  -\int_{0}^{t}D^{2}F\left(
s,{\phi}_{s}\left(  x\right)  \right)  D{\phi}_{s}\left(  x\right)  dW_{s}.
\end{align*}
From the properties of $F$  we readily have
\[
\int_{0}^{T}\left\|  \nabla F\left(  0, \cdot \right)  \right\|
_{L_{X}^{1} }dt<\infty,\quad P\text{-a.s.}
\]
%Moreover we have%
%\[
%\int_{0}^{T}\int_{B\left(  Y\right)  }\left|  \nabla F\left(  t,{y}\right)
%\right|  dydt<\infty,\quad P\text{-a.s.}%
%\]
%for every $P$-a.s. finite random variable $Y$, hence also%
%\[
%\int_{0}^{T}\int_{B\left(  Y\right)  }\left|  \nabla F\left(  t,{y}\right)
%\right|  J{\phi}_{t}^{-1}\left(  y\right)  dydt<\infty,\quad P\text{-a.s.}%
%\]
Because of the $P$-a.s. continuity of $D{\phi}_{t}\left(  x\right)
$ in $(  t,x)  $, we get
%\[
%\int_{0}^{T}\left\|  \nabla F\left(  t,{\phi}_{t}\left(  x\right)  \right)
%\right\|  _{L_{X}^{1}}dt<\infty,\quad P\text{-a.s.}%
%\]
%by a proper choice of $Y$ with respect to $X$ (using again the $P$-a.s.
%continuity of $D{\phi}_{t}\left(  x\right)  $ in $\left(  t,x\right)  $). For
%the same reason, this implies%
\[
\int_{0}^{T}\left\|  \nabla F\left(  t,{\phi}_{t}\left( \cdot
\right) \right) \cdot D{\phi}_{t}\left(  \cdot \right)  \right\|
_{L_{X}^{1}}dt<\infty,\quad
P\text{-a.s.}%
\]
which is another part of the estimate we need to prove.

Finally, we need to prove%
\[
\int_{0}^{T}\left\|  \int_{0}^{t}D^{2}F\left(  s,{\phi}_{s}\left(
\cdot\right)
\right)  D{\phi}_{s}\left(  \cdot\right)  dW_{s}\right\|  _{L_{X}^{1}}%
dt<\infty,\quad P\text{-a.s.}%
\]
If we prove that
\[
E\left[  \int_{0}^{T}\int_{K}\left|  \int_{0}^{t}D^{2}F\left(  s,{\phi}%
_{s}\left(  x\right)  \right)  D{\phi}_{s}\left(  x\right)  dW_{s}\right|
dxdt\right]  <\infty
\]
for every (deterministic) compact set $\mathbb{R} ^d$, then it is
not difficult to deduce the previous property;
% then the property follows. Up to a constant term
(using the inequality $\left| a\right|  \leq a^{2}+1$) it is
sufficient to bound
\begin{align*}
&  E\left[  \int_{0}^{T}\int_{K}\left|  \int_{0}^{t}D^{2}F\left(  s,{\phi}%
_{s}\left(  x\right)  \right)  D{\phi}_{s}\left(  x\right)  dW_{s}\right|
^{2}dxdt\right] \\
&  =\int_{0}^{T}\int_{K}E\left[  \left| \int_{0}^{t}D^{2}F\left(
s,{\phi }_{s}\left(  x\right)  \right) D{\phi}_{s}\left(  x\right)
dW_{s}\right|
^{2}\right]  dxdt\\
&  \leq\int_{0}^{T}\int_{K}E\left[  \int_{0}^{t}\left|
D^{2}F\left(  s,{\phi }_{s}\left(  x\right)  \right)
D{\phi}_{s}\left(  x\right)  \right|
^{2}ds\right]  dxdt\\
&  \leq T\int_{K}E\left[  \int_{0}^{T}\left|  D^{2}F\left(  s,{\phi}%
_{s}\left(  x\right)  \right)  D{\phi}_{s}\left(  x\right)  \right|
^{2}ds\right]  dx.
\end{align*}
 Let $p'$ such that $2 < p' < p$ and $2< p' < q$ where $p$ and $q$ are given
 in Hypothesis \ref{hy2}. Since
\[
\int_{K}\left(  \int_{0}^{T}E\left[  \left|  D{\phi}_{s}\left(
x\right) \right|  ^{r}\right]  ds\right)  dx<\infty,
\]
for every $r \ge 1$ (see (\ref{bound})), by the H\"older
inequality on $\Omega \times K \times [0,T]$,  it is sufficient to
prove that
\[
\int_{K}E\left[  \left(  \int_{0}^{T}\left|  D^{2}F\left(
s,{\phi}_{s}\left( x\right)  \right)  \right| ^{p'}ds\right)
\right] dx<\infty.
\]
But
%since
%$$
%\int_{\mathbb{R} ^d} f(\phi_s(x)) dx = \int_{\mathbb{R} ^d} \frac{f(y) dy}{ J \phi_s(\phi^{-1}_s(y))}
%$$
we have
\begin{align*}
E\left[  \left(  \int_{0}^{T}\int_{K}\left|  D^{2}F\left(
s,{\phi}_{s}\left( x\right)  \right)  \right| ^{p'}dxds\right)
\right]
\\
% =E\left[  \left( \int_{0}^{T}\int_{\phi_s (K)}\frac{\left|
%D^{2}F\left(  s,{y}\right) \right| ^{p'} dy}{ J \phi_s(\phi^{-1}_s(y))}ds\right) \right]
 =E\left[  \left( \int_{0}^{T}\int_{\phi_s (K)} \left|
D^{2}F\left(  s,{y}\right) \right| ^{p'}  J \phi^{-1}_s(y) dyds\right) \right]
\\
  \le \int_{0}^{T}\int_{\mathbb{R} ^d}\left|
D^{2}F\left(  s,{y}\right)  \right| ^{p'}E\left[
J{\phi}_{s}^{-1}\left(  y\right) \right] dyds.
\end{align*}
By the H\"older inequality,
$$
\int_{0}^{T}\int_{\mathbb{R} ^d}\left| D^{2}F\left(  s,{y}\right)
\right| ^{p'}E\left[ J{\phi}_{s}^{-1}\left(  y\right) \right] dyds
$$
$$
\le \int_{0}^{T} \Big ( \int_{\mathbb{R} ^d}\left| D^{2}F\left(
s,{y}\right) \right| ^{p}dy \Big )^{p'/p} \,
 \Big ( \int_{\mathbb{R} ^d}  (E\left| J{\phi}_{s}^{-1}\left(
y\right) \right |)^{\frac{p'}{p' -p}}dy \Big )^{\frac{p'-p}{p'}}
\, ds.
$$
Using the bound~(\ref{bound-inv}) we have
\[
 \sup_{x \in \mathbb{R} ^d}  \sup_{s \in [0,T]}  E \left[ \left| J{\phi}_{s}^{-1}\left(
x\right) \right|^{r}\right]<\infty
\]
 it is sufficient to
have
$$
\int_{0}^{T} \Big ( \int_{\mathbb{R} ^d}\left| D^{2}F\left(
s,{y}\right) \right| ^{p}dy \Big )^{p'/p} ds < \infty.
$$
This is true since $p'/p < q/p $ and
 $D^2 F
  \in L^q ([0,T]; L^p(\mathbb{R} ^d))$ by    Theorem \ref{KR}.
The proof is complete.
\end{proof}

\section{Counterexamples to certain extensions}

  \label{sec:counter} It would be interesting to try to generalize the results of this paper to random
vectorfields $b( x,t,\omega)$, possibly adapted. However our approach faces a fundamental difficulty:  it is very easy to exibith
a counterexample which shows that in some cases the regularizing effect disappears. Consider in one dimension the case
\[
b\left(  x,t,\omega\right)  =\sqrt{\left|  x-W_{t}\left(  \omega\right)
\right|  }%
\]
namely the stochastic differential equation
\[
dX_{t}^{x}=b\left(  X_{t}^{x},t,.\right)  dt+dW_{t},\quad t\geq0,\quad
X_{0}^{x}=x
\]
If $\left(  X_{t}^{x}\right)  $ is a solution, then $Y_{t}^{x}=X_{t}^{x}%
-W_{t}$ solves
\[
dY_{t}^{x}=\sqrt{\left|  Y_{t}^{x}\right|  }dt\quad t\geq0,\quad Y_{0}^{x}=x
\]
and viceversa. Hence the non-uniqueness for the latter equation transfer to
the former. In terms of stochastic transport equation, an equation of the
form
\[
\partial_{t}u\left(  t,x\right)  +\left(  b_{0}\left(  x-W\left(  t\right)
\right)  \cdot\nabla u\left(  t,x\right)  \right)  dt+\nabla u\left(
t,x\right)  \circ dW\left(  t\right)  =0
\]
may have several pathologies if $b_{0}$ is only H\"{o}lder continuous.

Unfortunately the previous example is not so artificial: something similar
happens in the nonlinear case. Let us argue only formally. Consider the
example in $\mathbb{R}^{2}$%
\[
\partial_{t}\xi\left(  t,x\right)  +\left(  u\left(  t,x\right)  \cdot
\nabla\xi\left(  t,x\right)  \right)  dt+\nabla\xi\left(  t,x\right)  \circ
dW\left(  t\right)  =0
\]
where $\xi=\partial_{2}u_{1}-\partial_{1}u_{2} $. This is the vorticity
equation of a 2D ideal fluid described by a stochastic version of Euler
equation. Following \cite{MajBar}, this equation is (formally) equivalent to
the stochastic ordinary equation depending on a parameter $a\in\mathbb{R}^{2}$%
\[
dX_{t}^{a}=\left[  \int_{\mathbb{R}^{2}}K( X_{t}^{a}-X_{t}^{a^{\prime}})
\xi_{0}( X_{t}^{a^{\prime}} ) da^{\prime}\right]  dt+dW_{t}%
\]
for a suitable kernel $K$, $\xi_{0}$ being the initial condition of the
vorticity equation. This problem is equivalent to
\[
dY_{t}^{a}=\left[  \int_{\mathbb{R}^{2}}K( Y_{t}^{a}-Y_{t}^{a^{\prime}})
\xi_{0}( X_{t}^{a^{\prime}}) da^{\prime}\right]  dt
\]
by the change of variable $Y_{t}^{a}=X_{t}^{a}-W_{t}$, and the equation for
$\left(  Y_{t}^{a}\right)  $ corresponds to the classical vorticity equation
\begin{align*}
\frac{\partial_{t}\xi^{\prime}\left(  t,x\right)  }{\partial t}+\left(
u^{\prime}\left(  t,x\right)  \cdot\nabla\xi^{\prime}\left(  t,x\right)
\right)  dt =0 \qquad\qquad\xi^{\prime} =\partial_{2}u_{1}^{\prime}%
-\partial_{1}u_{2}^{\prime}%
\end{align*}
with initial condition $\xi_{0}$. This means that the stochastic vorticity
equation is (at least formally) equivalent to the deterministic one. There is
no advantage to introduce that kind of stochastic perturbation.

%The case could
%be different for more complex stochastic perturbations for which the present analysis does not apply in a straighforward manner.



\bigskip

\bigskip

\begin{thebibliography}{99}
\bibitem{Ambrosio}Ambrosio, L., Transport equation and Cauchy problem for
$BV$ vector fields, Invent. Math. \textbf{158} (2004), no.~2, 227-260.

\bibitem {AmbrCrippa}Ambrosio, L., Crippa G., Existence, uniqueness,
stability and differentiability properties of the flow associated to weakly
differentiable vector fields, UMI-Lecture Notes, to appear.

%\bibitem {B}Bismut J. M., Martingales, the Malliavin Calculus and
%hypoellipticity general H\"{o}rmander's condition, Z.
%Wharscheinlichkeitstheorie Gabiele 56 (1981), 469-505.

\bibitem {Br}Brandt A., Interior Schauder estimates for parabolic differential
(or difference) equations via the maximum principle, Israel J. Math., Vol. 7
(1969), 254-262.

%\bibitem {Ce}Cerrai, S.: \textit{Second order PDE's in finite and infinite
%dimensions. A probabilistic approach.}, { Lectures Notes in Math.} 1762,
%Springer Verlag, 2001.

\bibitem {D}Davie A. M., Uniqueness of solutions of stochatic differential
equations, preprint (2007) (\texttt{arXiv:0709.4147}).

\bibitem {DiPernaLions}DiPerna, R. J.\ and\ Lions, P.-L., Ordinary
differential equations, transport theory and Sobolev spaces,
Invent. Math. \textbf{98} (1989), no.~3, 511-547.



\bibitem {Figalli} Figalli A., Existence and uniqueness of martingale solutions
for SDEs with rough or degenerate coefficients, 2007, to appear on \textit{J.
Funct. Anal.}

%\bibitem {F}Fourni\'{e}, E., Lasry, J.M., Lebuchoux, J., Lions, P.L. and
%Touzi, N., \textit{Applications of Malliavin calculus to Monte Carlo methods
%in finance,} Finance Stoch. 3 (1999), 391-412.

\bibitem {Kr}Krylov, N. V., {\it
 Lectures on elliptic and parabolic equations in
H\"{o}lder spaces,} Graduate Studies in Mathematics, 12. American
Mathematical Society, Providence, RI, 1996.

\bibitem{Kr1}  Krylov N. V., The heat equation in $L_q((0, T),L_p)-$spaces with
weights, SIAM J. on Math. Anal., Vol. 32, No. 5 (2001), 1117-1141.

%\bibitem {Kry2}{ Krylov N. V., } { Parabolic equations in $L_{p}$-spaces with
% mixed norms\/}, Algebra i Analiz., Vol. 14 (2002) , No. 4, 91-106 in Russian;
% English translation in St. Petersburg Math. J., Vol. 14 (2003), No. 4, 603-614.




\bibitem {KR}Krylov, N. V.; Rozovskii B. L., Stochastic partial differential
equations and diffusion processes. (Russian) Uspekhi Mat. Nauk 37 (1982), no.
6 (228), 75.95.

\bibitem {Kry-Ro}Krylov, N. V.; R\"{o}ckner, M., Strong solutions of stochastic
equations with singular time dependent drift. Probab. Theory Related Fields
131 (2005), no. 2, 154--196.


\bibitem{KP}Krylov, N. V.;  Priola, E., Elliptic and parabolic
 second-order PDEs with growing coefficients,  Preprint
  arXiv.org, June
 2008   (http://arxiv.org/abs/0806.3100v1).


\bibitem {K}Kunita, H., Stochastic differential equations and stochastic flows
of diffeomorphisms. Ecole d'\'{e}t\'{e} de probabilit\'{e}s de Saint-Flour,
XII---1982, 143--303, Lecture Notes in Math., 1097, Springer, Berlin, 1984.

\bibitem {K1}Kunita, H., First order stochastic partial differential
equations. Stochastic analysis (Katata/Kyoto, 1982), 249-269,
North-Holland Math. Library, 32, North-Holland, Amsterdam, 1984.

\bibitem {IW}Ikeda, N. and Watanabe, S., { Stochastic Differential Equations
and Diffusion Processes,} North Holland-Kodansha, 1981.

\bibitem {Lions}P. L. Lions, \textit{Mathematical Topics in Fluid Mechanics.
Volume 1: Incompressible models}, Oxford Univ. Press, Oxford, 1996.

%\bibitem {L}Lunardi, A.: \textit{Analytic semigroups and Optimal Regularity in
% Parabolic Problems, } Birkh\"{a}user, 1995.

\bibitem {MajBar}Majda, A. J. \ and\ Bertozzi, A. L.,
  \textit{Vorticity and
incompressible flow}, Cambridge Univ. Press, Cambridge, 2002.
%MR1867882 (2003a:76002)
%

%\bibitem{Pr} Priola E. : {\it Schauder estimates for a class of degenerate
% Kolmogorov equations}, Preprint 2008.
%

\bibitem{protter}
 Protter, P. E., \textit{ Stochastic Integration and Differential
Equations, Second Edition,} Springer, 2004.


\bibitem {V}Veretennikov, Yu. A., On strong solution and explicit formulas for
solutions of stochastic integral equations. Math. USSR Sb. 39, 387?403 (1981).
\end{thebibliography}
\end{document}
